\documentclass[aps,prx,preprint,amsmath,amssymb]{revtex4-2}
%\usepackage[bookmarks=false,linkcolor=NavyBlue,urlcolor=NavyBlue,colorlinks,citecolor=NavyBlue]{hyperref}
\usepackage[utf8]{inputenc}
\usepackage{natbib}
\usepackage{graphicx}
\usepackage{xcolor}
\usepackage{amsmath}
\usepackage{multirow}
\usepackage{array}
\usepackage{booktabs}
\usepackage[colorlinks=true, linkcolor=blue]{hyperref}
%\hypersetup{colorlinks,allcolors=black}
%\usepackage{amssymb}
%\usepackage{gensymb}
\usepackage{comment}
\usepackage{float}
%\usepackage{amsmath}
\usepackage{tabularx,graphicx}
\usepackage{epstopdf}
\usepackage{latexsym}
\usepackage{color, colortbl}
\usepackage{psfrag}
\usepackage{bbm}
\usepackage{bm}
\usepackage{titlesec}
\usepackage{dsfont}
\usepackage{feynmp}
\usepackage{slashed}
\usepackage{subcaption}
%\captionsetup[subfigure]{position=top, labelfont=bf,textfont=normalfont,singlelinecheck=off,justification=raggedright}
\usepackage{multirow}
%\usepackage{txfonts}
\usepackage{subcaption}
\usepackage[newcommands]{ragged2e}
\captionsetup{font=small, justification=Justified}
\usepackage[normalem]{ulem}
\renewcommand{\vec}[1]{{\boldsymbol{#1}}}
\usepackage{braket}
\usepackage{subcaption}
\def \rar{{\rightarrow}}
\def \uar{{\uparrow}}
\def \dar{{\downarrow}}
\def \k {{\vec k}}
\def \p {{\vec p}}
\def \R {{\cal{R}}}
\def \bI{{\vec I}}
\def \bJ{{\vec J}}
\def \e {\epsilon}
\def \ve {\varepsilon}
\def \r {{\vec r}}
\def \A {\mathrm{A}}
\def \B {\mathrm{B}}
\def \C {\mathrm{C}}
\def \v {{\bf v}}
\def \el {\ell}
\def \s {\vec{s}}
\def \q {{\vec q}}
\def \d{\partial}
\def \Q{{\cal Q}}
\def \Dg {\Delta_{\tn{gap}}}
\def \Ds {\Delta_{\tn{sc}}}
\def \l {\ell}
\def \dv {\delta\vec{\varphi}}
\def \S {{\cal{S}}}
\def \G {{\cal{G}}}
\def \C {{\cal{C}}}
\def \K {{\vec K}}
\def \x {{\vec x}}
\def \y {{\vec y}}
\def \Y {\mathbb{Y}}
\def \P {\vec{\Phi}}
\def \Ps {\Psi}
\def \Ph{\Phi}
\def \D{\Delta}
\def \t {\theta}
\def \ol{\overline}
\def \n{{\mathbf{n}}}
\def \L{{\cal{L}}}
\def \O {\cal{O}}
\def \beq {\begin{eqnarray}}
\def \eeq {\end{eqnarray}}
\def \tn {\textnormal}
\def \PP {{\cal {P}}}
\def \vp{\varphi}
\def \M{{\cal {M}}}
\def \Z {{\cal {Z}}}
\def \F {{\cal {F}}}
\def \E {{\cal {E}}}
\def \N {{\cal{N}}}
\def \m {{\mathcal}}
\def \nn {\nonumber}
\def \omd {{\omega_{\rm D}}}
\def \kpa {\kappa_\Vert}
\def \kpe {\kappa_\perp}
\def \ol{\overline}
\def \th{\vartheta}
\def \la{\langle}
\def \ra{\rangle}
\def \ua {\uparrow}
\def \da {\downarrow}
\def \cb {\color{blue}}
\def \ie {{\it i.e.}}
\def \phd {\phantomdagger}
\newcommand{\DC}[1]{ { \color{red} \footnotesize (\textsf{DC})\textsf{\textsl{#1}}}}
\newcommand{\HG}[1]{ { \color{green} \footnotesize (\textsf{HG})\textsf{\textsl{#1}}}}
\newcommand{\change}[1]{{\color{blue} #1}}
\newcommand{\rd}{{\rm d}}
\newcommand{\sgn}{{\rm sgn\,}}
\newcommand{\Tr}{{\rm \, Tr\,}}
\newcommand{\tr}{{\rm \, tr\,}}
\newcommand{\vn}[1]{{\left|\vec{#1}\right|}}
\newcommand{\calN}{{\mathcal N}}
\newcommand{\calG}{{\mathcal G}}
\newcommand{\calF}{{\mathcal F}}
\newcommand{\calD}{{\mathcal D}}
\newcommand{\calQ}{{\mathcal Q}}
\newcommand{\calL}{{\mathcal L}}
\newcommand{\calH}{{\mathcal H}}
\newcommand{\calX}{{\mathcal X}}
\newcommand{\calO}{{\mathcal O}}
\renewcommand{\Re}{\text{Re}}
\renewcommand{\Im}{\text{Im}}

\usepackage{tikz-feynman}
\usetikzlibrary{calc}
\usetikzlibrary{arrows}
\tikzset{
  % style to add an arrow in the middle of a path
  mid arrow/.style={postaction={decorate,decoration={
        markings,
        mark=at position .575 with {\arrow[#1]{stealth}}
      }}},
  near arrow/.style={postaction={decorate,decoration={
        markings,
        mark=at position .275 with {\arrow[#1]{stealth}}
      }}},
   far arrow/.style={postaction={decorate,decoration={
        markings,
        mark=at position .800 with {\arrow[#1]{stealth}}
      }}},
   boson/.style={decorate, draw=black,
    decoration={snake,amplitude=1pt, segment length=5pt},
      },
      phonon/.style={decorate, draw=red,
    decoration={snake,amplitude=1pt, segment length=5pt},
      },
   mid triangle/.style={postaction={decorate,decoration={
        markings,
        mark=at position .575 with {\arrow[#1]{triangle 45}}
      }}}
}


%\captionsetup{justification=justified}

\usepackage{tikz-feynman}
\usetikzlibrary{arrows}
\usetikzlibrary{arrows.meta, positioning,shapes.geometric}
\tikzset{
  % style to add an arrow in the middle of a path
  mid arrow/.style={postaction={decorate,decoration={
        markings,
        mark=at position .575 with {\arrow[#1]{stealth}}
      }}},
  near arrow/.style={postaction={decorate,decoration={
        markings,
        mark=at position .275 with {\arrow[#1]{stealth}}
      }}},
   far arrow/.style={postaction={decorate,decoration={
        markings,
        mark=at position .800 with {\arrow[#1]{stealth}}
      }}},
   boson/.style={decorate, draw=black,
    decoration={snake,amplitude=1pt, segment length=5pt},
      },
   mid triangle/.style={postaction={decorate,decoration={
        markings,
        mark=at position .575 with {\arrow[#1]{triangle 45}}
      }}}
}


\begin{document}
\title{Note on Monitored Chern Insulator (Reowrked)}

\author{Haoyu Guo}

\affiliation{Department of Physics, Cornell University, Ithaca, New York 14853, USA.}


\maketitle


\section{Bulk model}

We consider the following bulk model proposed by Matt: 
\begin{equation}\label{}
  h(k_x,k_y)/t_0=\sigma^1 \sin k_x +\sigma^2 \sin k_y+\sigma^3\left(r-\cos k_x-\cos k_y\right)\,.
\end{equation} This hamiltonian has a Chern number when $0<|r|<2$.

We can write down the real-space tight binding Hamiltonian as 
\begin{equation}\label{}
  h/t_0=\sigma^1\frac{T_x^\dagger-T_x}{2i}+\sigma^2 \frac{T_y^\dagger -T_y}{2i}+\sigma^3 \left(r-\frac{T_x+T_x^\dagger+T_y+T_y^\dagger}{2}\right)\,,
\end{equation} where $T_x,T_y$ are the translation operators in the $x,y$ directions respectively. 

\section{The model on the semi-infinite geometry}

Let's write down the second quantized hamiltonian 
\begin{equation}\label{}
\begin{split}
  H/t_{0}&= \sum_{k_x,y=0} c_{k_x,y}^\dagger\left(\sigma^1 \sin k_x+\sigma^3 (r-\cos k_x)\right)c_{kx,y} \\
&+ c_{k_x,y}^\dagger \left(\frac{\sigma^2}{2i}-\frac{\sigma^3}{2}\right)c_{k_x,y+1}+c_{k_x,y+1}^\dagger\left(-\frac{\sigma^2}{2i}-\frac{\sigma^3}{2}\right)c_{k_x,y}\,,
\end{split}
\end{equation} 

The Green's function is found to be 
\begin{equation}\label{}
  G^{-1}(z,k_x)=z-\sigma^1 \sin k_x- \sigma^3(r-\cos k_x)-\Sigma(z,k_x)\frac{1-\sigma^1}{2}\,,
\end{equation} where 
\begin{equation}\label{}
  \Sigma(z,k_x)=\frac{z^2+2r\cos k_x-r^2-\sqrt{z^2-E_+^2}\sqrt{z^2-E_-^2}}{2(z-\sin k_x)}\,.
\end{equation} and 
\begin{equation}\label{}
   E_{\pm}^2= \sin^2 k_x+(r-\cos k_x \pm 1)^2\,.
\end{equation} Here the branch cut is the standard one for square root.

The spectral function can be written as the sum of two terms $A=-2\Im G_R=A_\text{edge}+A_\text{bulk}$.

\begin{equation}\label{}
  A_\text{edge}=2\pi \delta(z-\sin k_x) \frac{1-(r-\cos k_x)^2}{2}\theta(1-(r-\cos k_x)^2)\begin{pmatrix}
           1 & 1 \\
           1 & 1 
         \end{pmatrix}\,,
\end{equation}
\begin{equation}\label{}
\begin{split}
  A_\text{bulk}&=\frac{1}{2(r-\cos k_x)^2}\frac{1}{|z-\sin k_x|}\sqrt{(E_+^2-z^2)(z^2-E_-^2)}\theta(E_-<|z|<E_+) \\
  &\times \begin{pmatrix}
                                                                                (r+z-\cos k_x-\sin k_x)^2 & (r-\cos k_x)^2-(z-\sin k_x)^2 \\
                                                                                 (r-\cos k_x)^2-(z-\sin k_x)^2 & (r-z-\cos k_x+\sin k_x)^2 
                                                                              \end{pmatrix}\,.
\end{split}
\end{equation}

\section{Setting up measurement}

We briefly discuss how the measurement is setup. The measurement is described by a set of Kraus operators:
\begin{equation}\label{}
  \rho\to K_\upsilon \rho K^\dagger_\upsilon\,,
\end{equation} with the POVM condition 
\begin{equation}\label{}
  \sum_{\upsilon} K_\upsilon^\dagger K_\upsilon = \mathds{1}\,.
\end{equation} We consider measuring a fermion operator
\begin{equation}
\hat{n}_i[\alpha]=(\alpha^\dagger c_i)^\dagger (\alpha^\dagger c_i)\,,
\end{equation} where $\alpha$ is a unit vector in the orbital space, so that $\hat{n}_i[\alpha]^2=\hat{n}_i[\alpha]$, i.e. it has eigenvalue 0,1.  Using this property, we can construct the Kraus operator using a $\upsilon_i^\alpha$ drawn from the real line. 
\begin{equation}\label{}
  K_{\upsilon_i^\alpha} \propto \exp\left(\kappa \upsilon_i^\alpha (\hat{n}_i[\alpha]-1/2) \right)\,,\quad \int \rd \upsilon_i^\alpha K_{\upsilon_i^\alpha}^\dagger K_{\upsilon_i^\alpha} \propto \mathds{1}\,.
\end{equation}  

Assuming continuum limit makes sense, the Keldysh action for the measurement takes the form 
\begin{equation}\label{}
  iS_m=\int\rd_t \rd^d\vec{r} \sum_{\alpha} \upsilon^\alpha(t,\vec{r})\sum_{a=1}^{r} \psi_a^\dagger \alpha \tau^x \alpha^\dagger \psi_a
\end{equation} 

Averaging over the trajectories with $\braket{\upsilon^\alpha \upsilon^\beta}=\gamma^2 \delta^{\alpha\beta}$, we therefore obtain 
\begin{equation}\label{}
  iS_m \Rightarrow \frac{\gamma^2}{2}\int\rd t\rd \vec{r} \sum_{\alpha} \psi^\dagger \mathbb{P}_\alpha \tau^x \psi  \mathbb{P}_\alpha \tau^x \psi\,.
\end{equation} Here $\mathbb{P}_\alpha=\alpha \alpha^\dagger$ is a projector in the orbital space. Decouple with Green's function $G\propto -i \psi \psi^\dagger$, we obtain 
\begin{equation}\label{}
  iS_m \Rightarrow \frac{\gamma^2}{2} \int_{t,\vec{r}} \sum_\alpha\Tr(G \tau^x \mathbb{P}_\alpha G \tau^x \mathbb{P}_\alpha)-f_R \Tr^2\left(G \tau^x \mathbb{P}_\alpha\right)
\end{equation} The trace is over replica, Keldysh, and orbital indices. To make sure the different Kraus operators commute, we require $\{\mathbb{P}_\alpha\}$ to be orthogonal to each other. For now we ignore the $f_R$ parameter, and just focus on the entanglement sector.

The full action is 
\begin{equation}\label{}
  iS = \ln\det\left(G_0^{-1}-\Sigma\right) +\Tr\left(\Sigma\cdot G\right)+i S_m
\end{equation}

We should note that while the saddle point is not very sensitive to the choice of $\{\mathbb{P}_\alpha\}$, but the fluctuations are. In particular, the fluctuations are diagonal in the basis described by $\{\mathbb{P}_\alpha\}$.

\section{Saddle point}

The saddle point equation with $f_R=0$ is 
\begin{equation}\label{}
  \Sigma_M(t,\vec{r})=-\gamma^2 \sum_\alpha\int\frac{\rd \omega}{2\pi}\frac{\rd k}{2\pi} \mathbb{P}_\alpha \tau^x G(\omega,k) \tau^x \mathbb{P}_\alpha =\frac{-i\Gamma}{2} \tau_z \sum_\alpha \mathbb{P}_\alpha\,,
\end{equation} Here
\begin{equation}\label{}
  \Gamma=\gamma^2\,,
\end{equation} and we have evaluated the $\omega$-integral using the sum rule (leading to a result proportional identity), and assume the BZ spans $k\in[-\pi,\pi]$.

\section{Orbit monitoring}

  Now we consider the orbit monitoring case, where $\{\mathbb{P}_\alpha\}=\{\mathbb{P}_{z,-},\mathbb{P}_{z,+}\}$, i.e. the eigenbasis of $\sigma_z$. In this setting, the saddle point self-energy is 
\begin{equation}\label{}
  \Sigma_{M,\text{saddle}}=-\frac{i \Gamma}{2}\tau_z\,.
\end{equation}

  We next derive the NLSM for the fluctuations around the saddle point. We first integrate out $G$, so the action becomes 
\begin{equation}\label{}
  i S = \ln \det\left(G_0^{-1}-\Sigma\right)+\frac{1}{2}\Tr\left(\Sigma\cdot  W_G^{-1}[\Sigma]\right)\,.
\end{equation} Here, the functional $W_G$ is local in space time, and the local action is given by 
\begin{equation}\label{}
  W_G[F]=-\gamma^2 \sum_{\alpha}\left[\tau^x \mathbb{P}_\alpha F \mathbb{P}_\alpha \tau^x-\tau^x \mathbb{P}_\alpha f_R \Tr\left(F \tau^x \mathbb{P}_\alpha\right)\right]
\end{equation}

\subsection{Choice of the parameter $f_R$}

  In this subsection, we determine the value of $f_R$ using the requirement that $\Sigma_M$ can be extended to a causal solution 
\begin{equation}\label{}
  \Sigma_M =\begin{pmatrix}
              \Sigma_R & \Sigma_K \\
              0 & \Sigma_A 
            \end{pmatrix}\,,
\end{equation} we assume all of them remain to be a frequency-independent number. Substituting the above ansatz into the saddle point equation, we obtain 
\begin{equation}\label{}
  \sum_{\alpha}\mathbb{P}_\alpha G_K \mathbb{P}_\alpha = \sum_{\alpha} f_R \mathbb{P}_\alpha \Tr\left(\mathbb{P}_\alpha G_K\right)\,.
\end{equation} Here $G_K$ is the local Keldysh Green's function of the saddle. The above equation simply says that the orbital-diagonal entries of $G_K$ is related to its trace over the replica indices, and therefore it requires 
\begin{equation}\label{}
  f_R=\frac{1}{R}\,,
\end{equation} which is the same conclusion as the previous AIII paper. The remaining determines the Keldysh self-energy $\Sigma_K$ in terms of the local Keldysh $G_K$.  In the upper-triangular Keldysh convention used throughout this note, $\Sigma_K$ is anti-Hermitian, so the Hermitian density kernel carries an additional factor of $-i$:
\begin{equation}\label{}
  \Sigma_K=-i \gamma^2\begin{pmatrix}
             G_K^1 & 0 \\
             0 & G_K^2 
           \end{pmatrix}_\text{orbit}=-i\Gamma\begin{pmatrix}
                                              1-2n_1 & 0 \\
                                              0 & 1-2n_2 
                                            \end{pmatrix}\,,
\end{equation} i.e. the Keldysh component is an anti-Hermitian matrix set by the orbital densities, while the Hermitian kernel is $\Gamma(1-2n_a)$. 

\subsection{Deriving the NLSM}

We again consider the half-filling case, so that the saddle has
\(\Sigma_K=0\).  Before choosing coordinates on the soft manifold, it is useful
to spell out the orbital structure of \(Q\).  The monitoring projectors are
\(\mathbb P_{z,+}\) and \(\mathbb P_{z,-}\), so the most general local
measurement saddle within the orbital-density Hubbard-Stratonovich channel is
block diagonal in this monitored orbital basis:
\begin{equation}
\label{eq:orbital-general-Q}
Q_{\rm orb}
=
\mathbb P_{z,+}\,Q_+
+\mathbb P_{z,-}\,Q_-,
\qquad
Q_\pm^2=\mathbf{1}.
\end{equation}
Here \(Q_\pm\) are matrices in Keldysh and replica space.  Equivalently,
\[
Q_\pm=T_\pm^{-1}\tau^zT_\pm
\]
with the usual causal constraints on \(T_\pm\).  In the explicit coordinates
used below one may write
\begin{equation}
\label{eq:orbital-general-Qpm}
Q_\pm
=
\cos\phi_\pm
\left[\tau^z\cos B_\pm+\tau^y\sin B_\pm\right]
+\sin\phi_\pm\,\tau^x,
\qquad
B_\pm=\theta_\pm\mathbf{1}_R+Y_\pm,
\end{equation}
with \(Y_\pm^\dagger=Y_\pm\) and \(\Tr_RY_\pm=0\).  Since the two orbital
projectors are orthogonal, \eqref{eq:orbital-general-Q} satisfies
\[
Q_{\rm orb}^2
=
\mathbb P_{z,+}Q_+^2+\mathbb P_{z,-}Q_-^2
=\mathbf{1}
\]
whenever \(Q_\pm^2=\mathbf{1}\).

Equation \eqref{eq:nlsm-Q-param} below is the diagonal submanifold
\[
Q_+=Q_-\equiv Q,
\qquad
Q_{\rm orb}=\mathbf{1}_{\rm orb}\,Q.
\]
This is the unique gapless orbital sector.  The reason is that only a common
Keldysh-replica rotation of both monitored orbitals is a symmetry of the full
fermion determinant.  A relative orbital rotation \(Q_+\neq Q_-\) inserts the
orbital matrix \(\mathbb P_{z,+}-\mathbb P_{z,-}=\sigma_z\) at the fluctuation
vertex.  The single-particle Hamiltonian contains orbital matrices that do not
commute with \(\sigma_z\), for example the \(\sin q\,\sigma_x\) term, so this
relative rotation is not protected by a Ward identity.  In the quadratic action
the common mode has the subtracted bubble
\(\Pi_{RA}(\Omega,k)-\Pi_{RA}(0,0)\), which vanishes at zero external arguments,
whereas the relative orbital mode involves bubbles with \(\sigma_z\) insertions,
such as
\[
\Pi_{RA}^{zz}(0,0)
=
\int\frac{d\omega\,dq}{(2\pi)^2}\,
\tr_{\rm orb}\!\left[
\sigma_zG_R(\omega,q)\sigma_zG_A(\omega,q)
\right],
\]
and this does not cancel the local measurement term.  Thus \(Q_+-Q_-\) is a
massive orbital-coherence mode.  Off-diagonal orbital components, which would
mix \(\mathbb P_{z,+}\) and \(\mathbb P_{z,-}\), are even more directly outside
the monitored-density saddle channel and are also massive.  The NLSM therefore
keeps only the orbital-identity component.

It is convenient to rewrite this remaining fluctuation manifold in a form that
makes the nonlinear constraint manifest.  Let
\begin{equation}\label{eq:nlsm-Q-param}
\Sigma=-\frac{i\Gamma}{2}Q,
\qquad
Q=\cos\phi\left[\tau^z\cos B+\tau^y\sin B\right]+\sin\phi\,\tau^x,
\end{equation}
where
\begin{equation}\label{eq:B-theta-Y}
B=\theta\,\mathbf{1}_R+Y,
\qquad
Y^\dagger=Y,
\qquad
\Tr_R Y=0.
\end{equation}
This is equivalent to the previous parametrization with
\(\mathcal X=e^{iB}=e^{i\theta}X\), \(X\in SU(R)\).  Since \(\cos B\) and
\(\sin B\) commute, one immediately checks that \(Q^2=\mathbf{1}\).  The saddle
point corresponds to
\begin{equation}\label{eq:Q-saddle}
Q_{\rm saddle}=\tau^z,
\qquad
\Sigma_{\rm saddle}=-\frac{i\Gamma}{2}\tau^z.
\end{equation}

For small fluctuations,
\begin{equation}\label{eq:Q-small-expand}
Q
=
\tau^z
+B\,\tau^y
+\phi\,\tau^x
-\frac{1}{2}\left(B^2+\phi^2\right)\tau^z
+O(\Phi^3),
\end{equation}
so the linear fluctuation of the self-energy is
\begin{equation}\label{eq:deltaSigma-linear}
\delta\Sigma
=
-\frac{i\Gamma}{2}\left(B\,\tau^y+\phi\,\tau^x\right)+O(\Phi^2).
\end{equation}

Expanding the action around the saddle gives
\begin{equation}\label{eq:quadratic-general}
iS
=
iS_{\rm saddle}
+\frac{1}{2}\Tr\left(\delta\Sigma\cdot W_G^{-1}[\delta\Sigma]\right)
-\frac{1}{2}\Tr\left(G_{\rm saddle}\,\delta\Sigma\,G_{\rm saddle}\,\delta\Sigma\right)
+O(\Phi^3).
\end{equation}
The linear term vanishes by the saddle-point equation.  Moreover, the
parametrization \eqref{eq:nlsm-Q-param} moves along the degenerate local saddle
manifold, so the purely local part of the quadratic kernel cancels between the
two terms in \eqref{eq:quadratic-general}.  The surviving quadratic action is
therefore controlled by the nonlocal fermion bubble.

We now Fourier transform the fluctuation fields and use the shorthand
\begin{equation}\label{eq:Q-int-short}
\int_Q \equiv \int \frac{d\Omega\,dk}{(2\pi)^2}.
\end{equation}
At linear order, the \(B\) mode couples through \(\tau^y\) and the \(\phi\) mode
through \(\tau^x\).  Using
\begin{equation}\label{eq:G-keldysh-triangular}
G(\omega,k)=
\begin{pmatrix}
G_R(\omega,k) & G_K(\omega,k)\\
0 & G_A(\omega,k)
\end{pmatrix},
\qquad
G_A=G_R^\dagger,
\end{equation}
the relevant Keldysh traces are
\begin{align}
\tr_K\!\left(G_1\tau^y G_2\tau^y\right)
&=
\tr_{\rm orb}\!\left(
G_{R,1}G_{A,2}+G_{A,1}G_{R,2}-G_{K,1}G_{K,2}
\right),
\label{eq:tauy-bubble}
\\
\tr_K\!\left(G_1\tau^x G_2\tau^x\right)
&=
\tr_{\rm orb}\!\left(
G_{R,1}G_{A,2}+G_{A,1}G_{R,2}+G_{K,1}G_{K,2}
\right),
\label{eq:taux-bubble}
\\
\tr_K\!\left(G_1\tau^x G_2\tau^y\right)
&=
i\,\tr_{\rm orb}\!\left(
G_{R,1}G_{A,2}-G_{A,1}G_{R,2}+G_{K,1}G_{K,2}
\right).
\label{eq:taux-tauy-mixing}
\end{align}
There is therefore a genuine mixed \(B\)-\(\phi\) quadratic kernel. Since the
saddle Green's function is proportional to the identity in replica space, the
\(U(1)\) and \(SU(R)\) parts of \(B\) share the same diagonal kernel, but only
the replica-singlet \(\theta\) mode can mix with \(\phi\):
\begin{equation}\label{eq:B-trace-split}
\Tr_R\!\left[B(-Q)B(Q)\right]
=
R\,\theta(-Q)\theta(Q)+\Tr_R\!\left[Y(-Q)Y(Q)\right].
\end{equation}
while
\begin{equation}\label{eq:B-phi-trace}
\Tr_R\!\left[B(-Q)\,\phi(Q)\right]=R\,\theta(-Q)\phi(Q),
\qquad
\Tr_R\!\left[Y(-Q)\,\phi(Q)\right]=0.
\end{equation}

Define the retarded-advanced bubble
\begin{equation}\label{eq:Pi-RA-def}
\Pi_{RA}(\Omega,k)
=
\int \frac{d\omega\,dq}{(2\pi)^2}\,
\tr_{\rm orb}\!\left[
G_R(\omega+\Omega,q+k)\,G_A(\omega,q)
\right],
\end{equation}
and the Keldysh-Keldysh bubble
\begin{equation}\label{eq:Pi-KK-def}
\Pi_{KK}(\Omega,k)
=
\int \frac{d\omega\,dq}{(2\pi)^2}\,
\tr_{\rm orb}\!\left[
G_K(\omega+\Omega,q+k)\,G_K(\omega,q)
\right].
\end{equation}
At the quadratic level, the two kernels entering the NLSM are
\begin{align}
\mathcal W_B(\Omega,k)
&=
\bigl[\Pi_{RA}(\Omega,k)+\Pi_{RA}(-\Omega,-k)-2\Pi_{RA}(0,0)\bigr]
\notag\\
&\qquad
-\bigl[\Pi_{KK}(\Omega,k)-\Pi_{KK}(0,0)\bigr],
\label{eq:WB-def}
\\
\mathcal W_\phi(\Omega,k)
&=
\bigl[\Pi_{RA}(\Omega,k)+\Pi_{RA}(-\Omega,-k)-2\Pi_{RA}(0,0)\bigr]
\notag\\
&\qquad
+\bigl[\Pi_{KK}(\Omega,k)-\Pi_{KK}(0,0)\bigr].
\label{eq:Wphi-def}
\end{align}
The off-diagonal kernels are
\begin{align}
\mathcal W_{\theta\phi}(\Omega,k)
&=
i\Bigl[
\Pi_{RA}(\Omega,k)-\Pi_{RA}(-\Omega,-k)
+\Pi_{KK}(\Omega,k)-\Pi_{KK}(0,0)
\Bigr],
\label{eq:Wthetaphi-def}
\\
\mathcal W_{\phi\theta}(\Omega,k)
&=
i\Bigl[
\Pi_{RA}(-\Omega,-k)-\Pi_{RA}(\Omega,k)
+\Pi_{KK}(\Omega,k)-\Pi_{KK}(0,0)
\Bigr].
\label{eq:Wphitheta-def}
\end{align}
Therefore
\begin{equation}\label{eq:nlsm-quadratic-final}
iS^{(2)}
=
-\frac{\Gamma^2}{8}\int_Q
\left[
\mathcal W_B(\Omega,k)\,
\Bigl(R\,\theta(-Q)\theta(Q)+\Tr_R[Y(-Q)Y(Q)]\Bigr)
+R\,\mathcal W_{\theta\phi}(\Omega,k)\,\theta(-Q)\phi(Q)
+R\,\mathcal W_{\phi\theta}(\Omega,k)\,\phi(-Q)\theta(Q)
+R\,\mathcal W_\phi(\Omega,k)\,\phi(-Q)\phi(Q)
\right].
\end{equation}

The retarded-advanced part is the same fermion bubble analyzed in
Appendix~\ref{app:ra-bubble}.  In the notation introduced there,
\begin{equation}\label{eq:Delta-to-W}
\Gamma\,\Re\!\left[
\frac{\Pi_{RA}(\Omega,k)+\Pi_{RA}(-\Omega,-k)}{2}-\Pi_{RA}(0,0)
\right]
=
\Delta(\Omega,k),
\end{equation}
so the real part of the retarded-advanced contribution to the NLSM kernel is
\begin{equation}\label{eq:W-RA-real}
\Re\,\mathcal W_{B,\phi}^{RA}(\Omega,k)=\frac{2}{\Gamma}\,\Delta(\Omega,k).
\end{equation}
For the off-diagonal piece, the retarded-advanced contribution is the odd part
\begin{equation}\label{eq:W-RA-mixed}
\mathcal W_{\theta\phi}^{RA}(\Omega,k)
=
i\Bigl[\Pi_{RA}(\Omega,k)-\Pi_{RA}(-\Omega,-k)\Bigr],
\end{equation}
so the \(\theta\)-\(\phi\) coupling starts at odd order in the external
variables.  Moreover,
\begin{equation}\label{eq:PiKK-even}
\Pi_{KK}(-\Omega,-k)=\Pi_{KK}(\Omega,k),
\end{equation}
which follows directly from the definition \eqref{eq:Pi-KK-def} by shifting the
integration variables and using cyclicity of the orbital trace.  Therefore the
Keldysh-Keldysh contribution to the mixed kernel is even in \((\Omega,k)\) and,
after subtracting the local piece \(\Pi_{KK}(0,0)\), it starts only at even
order.  In particular, the leading odd small-\((\Omega,k)\) term remains
proportional to \(\Omega k\) and comes entirely from the retarded-advanced
sector.
Hence, if we leave the Keldysh-Keldysh piece abstract for the moment, the
quadratic NLSM inherits the previously established small-\((\Omega,k)\) structure
\begin{equation}\label{eq:nlsm-small-structure}
\Delta(\Omega,k)
=
-A_k\,k^2
+A_{\Omega k}\,\Omega k
-\frac{\Omega^2}{8\pi}\log\!\frac{C_\Omega}{\Gamma|\Omega|}
+\cdots.
\end{equation}
The \(U(1)\) mode \(\theta\) and the \(SU(R)\) mode \(Y\) therefore share the
same singular retarded-advanced diagonal kernel, the \(\phi\) mode differs by
the sign of the abstract \(\Pi_{KK}\) contribution, and there is in general a
\(\theta\)-\(\phi\) mixing kernel \(\mathcal W_{\theta\phi}\) as well.

\subsubsection{Diagonalizing the \(\theta\)-\(\phi\) sector for \(\Pi_{KK}=0\)}

If we temporarily set
\begin{equation}
\Pi_{KK}(\Omega,k)=0,
\end{equation}
then the singlet \(\theta\)-\(\phi\) block of
\eqref{eq:nlsm-quadratic-final} takes the form
\begin{equation}\label{eq:theta-phi-block}
\begin{split}
iS^{(2)}_{\theta\phi}
=
-\frac{\Gamma^2R}{8}\int_Q
\Bigl[
W_s(\Omega,k)\,\theta(-Q)\theta(Q)
+W_s(\Omega,k)\,\phi(-Q)\phi(Q)
\\
\qquad\qquad\qquad\qquad
+W_a(\Omega,k)\,\theta(-Q)\phi(Q)
-W_a(\Omega,k)\,\phi(-Q)\theta(Q)
\Bigr],
\end{split}
\end{equation}
where
\begin{align}
W_s(\Omega,k)
&=
\Pi_{RA}(\Omega,k)+\Pi_{RA}(-\Omega,-k)-2\Pi_{RA}(0,0),
\label{eq:Ws-def}
\\
W_a(\Omega,k)
&=
i\Bigl[\Pi_{RA}(\Omega,k)-\Pi_{RA}(-\Omega,-k)\Bigr].
\label{eq:Wa-def}
\end{align}

Now define the circular combinations
\begin{equation}\label{eq:psi-pm-def}
\psi_\pm(Q)=\frac{\theta(Q)\pm i\phi(Q)}{\sqrt2},
\qquad
\bar\psi_\pm(Q)=\frac{\theta(-Q)\mp i\phi(-Q)}{\sqrt2}.
\end{equation}
In terms of these variables,
\begin{align}
\bar\psi_+(Q)\psi_+(Q)
&=
\frac{1}{2}\Bigl[
\theta(-Q)\theta(Q)+\phi(-Q)\phi(Q)
\notag\\
&\qquad\qquad\qquad
+i\theta(-Q)\phi(Q)-i\phi(-Q)\theta(Q)
\Bigr],
\\
\bar\psi_-(Q)\psi_-(Q)
&=
\frac{1}{2}\Bigl[
\theta(-Q)\theta(Q)+\phi(-Q)\phi(Q)
\notag\\
&\qquad\qquad\qquad
-i\theta(-Q)\phi(Q)+i\phi(-Q)\theta(Q)
\Bigr].
\end{align}
Therefore the singlet block diagonalizes as
\begin{equation}\label{eq:theta-phi-diagonalized}
iS^{(2)}_{\theta\phi}
=
-\frac{\Gamma^2R}{8}\int_Q
\left[
\bigl(W_s-iW_a\bigr)\,\bar\psi_+(Q)\psi_+(Q)
+\bigl(W_s+iW_a\bigr)\,\bar\psi_-(Q)\psi_-(Q)
\right].
\end{equation}
Using \eqref{eq:Ws-def} and \eqref{eq:Wa-def}, the two eigen-kernels are simply
\begin{align}
W_s(\Omega,k)-iW_a(\Omega,k)
&=
2\Bigl[\Pi_{RA}(\Omega,k)-\Pi_{RA}(0,0)\Bigr],
\\
W_s(\Omega,k)+iW_a(\Omega,k)
&=
2\Bigl[\Pi_{RA}(-\Omega,-k)-\Pi_{RA}(0,0)\Bigr].
\end{align}
So when \(\Pi_{KK}=0\), the \(\theta\)-\(\phi\) sector is exactly diagonalized by
the circular combinations \(\psi_\pm\): one mode is governed by the retarded-
advanced bubble at \((\Omega,k)\), and the other by the same bubble at
\((-\Omega,-k)\).  The mixing seen in the \((\theta,\phi)\) basis is therefore
just the statement that the natural eigenmodes are \(\theta\pm i\phi\), rather
than \(\theta\) and \(\phi\) separately.

\subsubsection{Including the even \(\Pi_{KK}\) contribution}

When \(\Pi_{KK}\neq 0\), the symmetry \eqref{eq:PiKK-even} implies that only the
even combination
\begin{equation}
K_s(\Omega,k)\equiv \Pi_{KK}(\Omega,k)-\Pi_{KK}(0,0)
\end{equation}
can enter the singlet block.  Then \eqref{eq:theta-phi-block} generalizes to
\begin{equation}
\begin{split}
iS^{(2)}_{\theta\phi}
=
-\frac{\Gamma^2R}{8}\int_Q
\Bigl[
W_s\,\theta(-Q)\theta(Q)
+W_s\,\phi(-Q)\phi(Q)
+W_a\,\theta(-Q)\phi(Q)
-W_a\,\phi(-Q)\theta(Q)
\\
\qquad\qquad\qquad\qquad
-K_s\,\theta(-Q)\theta(Q)
+K_s\,\phi(-Q)\phi(Q)
+iK_s\,\theta(-Q)\phi(Q)
+iK_s\,\phi(-Q)\theta(Q)
\Bigr].
\end{split}
\end{equation}
Transforming to the circular variables \(\psi_\pm\) gives
\begin{equation}\label{eq:theta-phi-with-KK}
iS^{(2)}_{\theta\phi}
=
-\frac{\Gamma^2R}{8}\int_Q
\left[
\bigl(W_s-iW_a\bigr)\,\bar\psi_+(Q)\psi_+(Q)
+\bigl(W_s+iW_a\bigr)\,\bar\psi_-(Q)\psi_-(Q)
-2K_s(\Omega,k)\,\bar\psi_+(Q)\psi_-(Q)
\right].
\end{equation}
Using \eqref{eq:Ws-def} and \eqref{eq:Wa-def}, this can be written as
\begin{equation}
\begin{split}
iS^{(2)}_{\theta\phi}
=
-\frac{\Gamma^2R}{4}\int_Q
\Bigl[
\bigl(\Pi_{RA}(\Omega,k)-\Pi_{RA}(0,0)\bigr)\bar\psi_+(Q)\psi_+(Q)
\\
\qquad\qquad
+\bigl(\Pi_{RA}(-\Omega,-k)-\Pi_{RA}(0,0)\bigr)\bar\psi_-(Q)\psi_-(Q)
-K_s(\Omega,k)\,\bar\psi_+(Q)\psi_-(Q)
\Bigr].
\end{split}
\end{equation}
So the effect of a nonzero \(\Pi_{KK}\) is not to change the odd
\(\Omega k\) coefficient of the diagonalized retarded-advanced kernels.  Instead,
it introduces an even off-diagonal mixing between the two circular modes
\(\psi_+\) and \(\psi_-\).  The simple diagonalization by \(\theta\pm i\phi\)
therefore holds only when \(\Pi_{KK}=0\); otherwise those combinations remain
useful because they isolate the retarded-advanced odd part, but they are coupled
to each other by the even kernel \(K_s(\Omega,k)\).

\subsubsection{Low-energy dispersion of the \texorpdfstring{$(\theta,\phi)$}{(theta,phi)} sector}

To identify the low-energy singlet modes, we now follow the same strategy as in
the Guo et al. reference in the \texttt{References/} folder: first diagonalize the \((\theta,\phi)\) block in the
circular basis, then insert the low-energy bubble expansion, and finally read
off the resulting soft trajectories and pole structure.

As discussed above, the full KK bubble is smooth at the origin and extremely
small numerically; Appendix~\ref{app:kk-bubble} gives
\[
\Re \Pi_{KK}(0,0)
\approx
-8.38\times 10^{-5},\quad
-8.08\times 10^{-6},\quad
-6.70\times 10^{-7}
\qquad
(\Gamma=20,40,80),
\]
consistent with the large-\(\Gamma\) estimate
\(\Pi_{KK}=O(\Gamma^{-4})\); see \eqref{eq:PiKK-largeGamma-appendix}.  If we
therefore treat \(\Pi_{KK}\) as independent of \((\Omega,k)\) at low energy,
then \(K_s(\Omega,k)=\Pi_{KK}(\Omega,k)-\Pi_{KK}(0,0)\approx 0\), and the
singlet block is controlled entirely by the retarded-advanced kernel.

Using \eqref{eq:theta-phi-diagonalized}, the quadratic action becomes
\begin{equation}\label{eq:psi-pm-lowE-general}
iS^{(2)}_{\theta\phi}
\simeq
-\frac{\Gamma^2R}{4}\int_Q
\Bigl[
\delta\Pi_{RA}(\Omega,k)\,\bar\psi_+(Q)\psi_+(Q)
\;+\;
\delta\Pi_{RA}(-\Omega,-k)\,\bar\psi_-(Q)\psi_-(Q)
\Bigr],
\end{equation}
where
\[
\delta\Pi_{RA}(\Omega,k)\equiv \Pi_{RA}(\Omega,k)-\Pi_{RA}(0,0).
\]
The appendix result \eqref{eq:Delta-summary-app} then implies, to leading order
in small external arguments,
\begin{equation}\label{eq:psi-pm-lowE-real}
\Gamma\,\Re\,\delta\Pi_{RA}(\pm\Omega,\pm k)
=
-A_k\,k^2
\pm A_{\Omega k}\,\Omega k
-\frac{\Omega^2}{8\pi}\log\!\frac{C_\Omega}{\Gamma|\Omega|}
+\cdots,
\end{equation}
with
\begin{equation}\label{eq:psi-pm-coeffs}
A_k\xrightarrow[\Gamma\to\infty]{}\frac{1}{3\sqrt3}-\frac18,
\qquad
A_{\Omega k}\xrightarrow[\Gamma\to\infty]{}\frac{1}{3\sqrt3}.
\end{equation}
The even real part is not the whole story, however.  The full complex bubble can
also contain terms linear in \(\Omega\) and \(k\).  As shown in
Appendix~\ref{app:ra-bubble-odd-linear}, those linear terms are necessarily
purely imaginary:
\begin{equation}\label{eq:psi-pm-linear-imag}
\Gamma\,\delta\Pi_{RA}(\Omega,k)
=
i\bigl[b_\Omega(\Gamma)\,\Omega-b_k(\Gamma)\,k\bigr]
-A_k\,k^2
+A_{\Omega k}\,\Omega k
-\frac{\Omega^2}{8\pi}\log\!\frac{C_\Omega}{\Gamma|\Omega|}
+\cdots,
\end{equation}
with
\begin{equation}\label{eq:odd-linear-scaling-main}
b_\Omega(\Gamma)=O(\Gamma^{-1}),
\qquad
b_k(\Gamma)=O(\Gamma^{-1}).
\end{equation}
These odd analytic terms therefore do not modify the real-part zero condition
used below to identify the soft real-frequency trajectories, although they do
contribute to the phase of the circular \(\psi_\pm\) kernels and hence to the
\(\theta\)-\(\phi\) mixing structure.
Since \(\Omega\) in \eqref{eq:Pi-RA-def} is already the physical real
frequency, there is no additional analytic continuation here.  The low-energy
dispersion is read off directly from the real-frequency quadratic kernels of
\(\psi_\pm\).

It is convenient to define the dimensionless ratio
\begin{equation}\label{eq:v-theta-phi-def}
v\equiv \frac{A_{\Omega k}}{A_k},
\qquad
v\xrightarrow[\Gamma\to\infty]{}
\frac{1/(3\sqrt3)}{1/(3\sqrt3)-1/8}
\approx 2.853.
\end{equation}
Then \eqref{eq:psi-pm-lowE-real} becomes
\begin{equation}\label{eq:psi-pm-lowE-real-v}
\Gamma\,\Re\,\delta\Pi_{RA}(\pm\Omega,\pm k)
=
-A_k
\left[
k^2\mp v\,\Omega k
+\frac{\Omega^2}{8\pi A_k}\log\!\frac{C_\Omega}{\Gamma|\Omega|}
\right]
+\cdots.
\end{equation}
If one temporarily drops the nonanalytic \(\Omega^2\log|\Omega|\) term, the
two circular modes would have the linear dispersions
\begin{equation}\label{eq:linear-dispersion-theta-phi}
\Omega_\pm(k)\simeq \pm \frac{k}{v},
\qquad
\frac{1}{v}\approx 0.350.
\end{equation}
This is the direct analogue of the shear structure found in the Guo et al.
analysis: the \(\psi_+\) and \(\psi_-\) sectors disperse in opposite
directions.

Keeping the logarithmic term, the zero of the real part of the quadratic kernel
obeys
\begin{equation}\label{eq:dispersion-equation-theta-phi}
k^2\mp v\,\Omega k
+\frac{\Omega^2}{8\pi A_k}\log\!\frac{C_\Omega}{\Gamma|\Omega|}
=0.
\end{equation}
For fixed \(\Omega\), this is a quadratic equation for \(k\):
\begin{equation}\label{eq:k-of-omega-theta-phi}
k_{\pm,\sigma}(\Omega)
=
\frac{\Omega}{2A_k}
\left[
\pm A_{\Omega k}
+\sigma\sqrt{
A_{\Omega k}^2
-\frac{A_k}{2\pi}\log\!\frac{C_\Omega}{\Gamma|\Omega|}
}
\right],
\qquad
\sigma=\pm.
\end{equation}
Real solutions exist only when
\begin{equation}\label{eq:discriminant-theta-phi}
A_{\Omega k}^2
>
\frac{A_k}{2\pi}\log\!\frac{C_\Omega}{\Gamma|\Omega|}.
\end{equation}
Because the logarithm grows without bound as \(|\Omega|\to 0\), this condition
eventually fails at asymptotically low frequency.  Therefore the
\((\theta,\phi)\) sector does not possess an exact real-frequency pole that
survives all the way to the origin.  Instead, the linear relation
\(\Omega_\pm(k)\simeq \pm k/v\) describes the incipient dispersion that is
progressively cut off by the singular \(\Omega^2\log(1/|\Omega|)\) term.

Even when the discriminant is negative, the kernel still has a soft ridge.
Completing the square in \eqref{eq:psi-pm-lowE-real-v} gives
\begin{equation}\label{eq:soft-ridge-theta-phi}
k^2\mp v\,\Omega k
+\frac{\Omega^2}{8\pi A_k}\log\!\frac{C_\Omega}{\Gamma|\Omega|}
=
\left(k\mp\frac{v}{2}\Omega\right)^2
+\left[
\frac{1}{8\pi A_k}\log\!\frac{C_\Omega}{\Gamma|\Omega|}
-\frac{v^2}{4}
\right]\Omega^2.
\end{equation}
Thus the minimum of the real-frequency kernel at fixed \(\Omega\) lies along
\begin{equation}\label{eq:ridge-theta-phi}
k_\pm^{\ast}(\Omega)=\pm \frac{v}{2}\,\Omega,
\end{equation}
while the residual stiffness along this ridge is set by the bracketed
\(\Omega^2\) coefficient in \eqref{eq:soft-ridge-theta-phi}.  At sufficiently
small \(|\Omega|\), the logarithm dominates, so the ridge remains soft but does
not turn into a true zero mode away from \((\Omega,k)=(0,0)\).

\subsection{Implications of the bubble analysis}

The technical analysis of the fermion bubble is collected in
Appendix~\ref{app:bubble-appendix}.  Here we only summarize the conclusions that
enter the NLSM.

First, the retarded-advanced bubble controls the singular low-energy structure.
In the notation of \eqref{eq:Delta-to-W},
\begin{equation}\label{eq:Delta-summary-main}
\Delta(\Omega,k)
=
-A_k\,k^2
+A_{\Omega k}\,\Omega k
-\frac{\Omega^2}{8\pi}\log\!\frac{C_\Omega}{\Gamma|\Omega|}
+\cdots,
\end{equation}
with
\begin{equation}\label{eq:Ak-Aomegak-summary-main}
A_k\xrightarrow[\Gamma\to\infty]{}\frac{1}{3\sqrt{3}}-\frac{1}{8},
\qquad
A_{\Omega k}\xrightarrow[\Gamma\to\infty]{}\frac{1}{3\sqrt{3}}.
\end{equation}
These coefficients come from the exact \(--\) sector of the full Green's
function rather than from the bare edge pole alone; see
Appendix~\ref{app:ra-bubble}.

Second, the Keldysh-Keldysh bubble is even under
\((\Omega,k)\rightarrow(-\Omega,-k)\) and therefore does not affect the leading
odd \(\Omega k\) term.  More importantly, the full KK bubble is smooth near the
origin and is strongly suppressed at large \(\Gamma\); see
Appendix~\ref{app:kk-bubble}.  Thus the nonanalytic low-energy structure of the
quadratic NLSM kernel comes entirely from the retarded-advanced sector.

\subsection{Replica sector \(Y\)}

The traceless replica field \(Y\) is simpler than the singlet sector because it
does not mix with \(\phi\).  From \eqref{eq:nlsm-quadratic-final}, its quadratic
action is
\begin{equation}\label{eq:Y-sector-general}
iS_Y^{(2)}
=
-\frac{\Gamma^2}{8}\int_Q \mathcal W_Y(\Omega,k)\,
\Tr_R\!\left[Y(-Q)Y(Q)\right],
\end{equation}
with
\begin{equation}\label{eq:Y-kernel-general}
\mathcal W_Y(\Omega,k)\equiv \mathcal W_B(\Omega,k)
=
\bigl[\Pi_{RA}(\Omega,k)+\Pi_{RA}(-\Omega,-k)-2\Pi_{RA}(0,0)\bigr]
-\bigl[\Pi_{KK}(\Omega,k)-\Pi_{KK}(0,0)\bigr].
\end{equation}
So the \(Y\) kernel is exactly the same diagonal \(\tau^y\) kernel as the
replica-singlet field \(\theta\), but without any coupling to \(\phi\).

Because the retarded-advanced bubble carries the singular
small-\((\Omega,k)\) dependence while the full KK bubble is smooth and strongly
suppressed, the leading low-energy behavior of the \(Y\) sector is the same as
that of the \(\theta\) sector.  In particular, the dominant momentum, mixed, and
frequency dependence are inherited from \eqref{eq:Delta-summary-main}.

\section{Generic-\texorpdfstring{$r$}{r} Fermion-Bubble Coefficients}
\label{sec:generic-r-bubbles}

The preceding fermion-bubble analysis was written at the symmetric value
\(r=1\).  Here we keep the same Chern-insulator phase,
\begin{equation}
0<r<2,
\end{equation}
and redo the large-\(\Gamma\) coefficient extraction.  To avoid overloading the
same symbol in several roles, this section uses the following conventions:
\(r\) is the microscopic mass parameter of the Chern-insulator Hamiltonian,
\(q_r\) is the positive momentum at which the edge mode touches the bulk
continuum, \(c_r\) is the edge velocity at that touching point, and \(s_r\) is
the corresponding sine.  Explicitly,
\begin{equation}
q_r\equiv \arccos(r-1),
\qquad
c_r\equiv r-1,
\qquad
s_r\equiv \sin q_r=\sqrt{r(2-r)}.
\end{equation}
Thus \(c_r\) is not a new coupling: it is simply
\(\cos q_r\), or equivalently the endpoint value of the group velocity
\(\partial_q\sin q\).  The main change relative to \(r=1\) is that the
edge-bulk touching points move from \(\pm\pi/2\) to \(\pm q_r\).
The edge pole exists on the interval
\begin{equation}
\mathcal I_r=(-q_r,q_r),
\qquad
|r-\cos q|<1,
\end{equation}
and its residue is
\begin{equation}
Z_r(q)\equiv 1-(r-\cos q)^2.
\end{equation}
Inside \(\mathcal I_r\), the bulk self-energy has the pole
\begin{equation}
\Sigma_R(\omega,q)
=
\frac{Z_r(q)}{\omega-\sin q+i0}
\;+\;
\Sigma_{\rm reg}(\omega,q),
\end{equation}
while outside \(\mathcal I_r\) this pole is absent.  In the \(\sigma_x\)
eigenbasis, the retarded inverse Green's function is
\begin{equation}
G_R^{-1}(\omega,q)
=
\begin{pmatrix}
\omega-\sin q+i\Gamma/2 & -(r-\cos q)\\
-(r-\cos q) & \omega+\sin q-\Sigma_R(\omega,q)+i\Gamma/2
\end{pmatrix}.
\end{equation}
As before, the large-\(\Gamma\) small-argument coefficients are dominated by the
\(--\) component.  On the resonance strip
\begin{equation}
\omega-\sin q=O\!\left(\frac{Z_r(q)}{\Gamma}\right),
\end{equation}
the leading local form is
\begin{equation}
\label{eq:Gmm-local-generic-r}
G_{--}^{\rm loc}(\omega,q)
=
\frac{\omega-\sin q}
{i\Gamma(\omega-\sin q)/2-Z_r(q)}.
\end{equation}

For the reduced local bubble, the \(\omega\) integral is again elementary.  After
taking the real part and subtracting the local \((\Omega,k)=(0,0)\) piece, the
small-\((\Omega,k)\) density is
\begin{equation}
\label{eq:generic-r-density}
\Gamma\,\Re\,\mathcal I_\Gamma(q;\Omega,k)
=
-\frac{4Z_r(q)}{\Gamma^2}
-\frac{(\Omega-k\cos q)^2}{4Z_r(q)}
+O\!\left(Q^3,\frac{Q^2}{\Gamma^2},\frac{k}{\Gamma^2}\right),
\end{equation}
where \(Q\) denotes either external argument.  Therefore the formal local Taylor
kernel is
\begin{equation}
\label{eq:generic-r-formal-kernel}
\Delta_r^{(--)}(\Omega,k)
=
-\int_{-q_r}^{q_r}\frac{dq}{8\pi}\,
\frac{(\Omega-k\cos q)^2}{Z_r(q)}
+\cdots.
\end{equation}
At \(r=1\), the factors of \(\cos q\) in the \(k^2\) and \(\Omega k\) terms
remove the endpoint divergence, giving the finite constants quoted above.  For a
generic \(r\neq1\), however, the endpoint velocity is
\(\cos q_r=c_r\neq0\), so the same endpoint that produced the
\(\Omega^2\log|\Omega|\) term now also produces logarithms in the momentum and
mixed coefficients.

Near either touching point,
\begin{equation}
Z_r(q_r-p)=2s_r p+O(p^2),
\qquad
Z_r(-q_r+p)=2s_r p+O(p^2),
\end{equation}
and \(\cos q\to c_r\).  Thus the singular endpoint contribution is
\begin{equation}
\label{eq:generic-r-singular-kernel}
\Delta_{r,\rm sing}^{(--)}(\Omega,k)
=
-\frac{(\Omega-c_r k)^2}{8\pi s_r}
\log\!\frac{C_r}{\Gamma|\Omega-c_r k|}
+O(Q^2),
\end{equation}
where \(Q\) denotes a generic small external argument, either \(\Omega\) or
\(k\).  The combination
\(\Omega-c_r k\) is the mismatch between the external frequency and the edge
dispersion linearized at either touching point in the bubble convention used
above:
\[
\Omega+\sin(q-k)-\sin q \simeq \Omega-c_r k
\qquad
(q\simeq \pm q_r).
\]
The constant \(C_r\) is a nonuniversal matching constant.  It absorbs the
finite part of the pure endpoint integral proportional to
\((\Omega-c_r k)^2\), and it also depends on the precise convention used to
match the local endpoint model onto the nonsingular part of the full lattice
Green's function.  The important universal statement in
\eqref{eq:generic-r-singular-kernel} is therefore the coefficient of the
logarithm.

The infrared cutoff in this logarithm is fixed by the point where the
fixed-\(q\) Taylor expansion \eqref{eq:generic-r-formal-kernel} stops being
valid.  Near the right touching point, \(q=q_r-p\), the residue is
\[
Z_r(q)\simeq 2s_r p.
\]
The unexpanded local kernel contains the denominator
\[
\bigl[Z_r(q)+Z_r(q-k)\bigr]^2
+\frac{\Gamma^2}{4}
\bigl[\Omega+\sin(q-k)-\sin q\bigr]^2.
\]
Thus the Taylor expansion in external arguments assumes
\[
Z_r(q)\gg \Gamma|\Omega-c_r k|
\]
up to order-one constants.  It fails when
\[
p\lesssim p_{\min},
\qquad
p_{\min}\sim \frac{\Gamma|\Omega-c_r k|}{2s_r}.
\]
The endpoint integral therefore has the structure
\[
\int_{p_{\min}}^{p_{\rm UV}}\frac{dp}{p}
=
\log\frac{p_{\rm UV}}{p_{\min}}
=
\log\frac{C_r}{\Gamma|\Omega-c_r k|},
\]
where the upper matching scale and all order-one constants have been absorbed
into \(C_r\).  When \(\Omega\simeq c_r k\), the leading endpoint mismatch is
tuned away and subleading terms in the external expansion regulate the kernel;
along that direction one should use the unexpanded endpoint kernel rather than
only \eqref{eq:generic-r-singular-kernel}.

With this convention for \(C_r\), the remaining regular endpoint term is
unambiguous and comes from the convergent difference
\begin{equation}
\label{eq:generic-r-regular-integral}
\Delta_{r,\rm reg}^{(--)}(\Omega,k)
=
-\frac{1}{8\pi}
\int_{-q_r}^{q_r}dq\,
\frac{(\Omega-k\cos q)^2-(\Omega-c_r k)^2}{Z_r(q)}.
\end{equation}
Using
\[
Z_r(q)=(\cos q-c_r)(r+1-\cos q),
\]
the numerator of the subtracted expression factorizes as
\[
(\Omega-k\cos q)^2-(\Omega-c_r k)^2
=
(\cos q-c_r)\left[-2\Omega k+k^2(\cos q+c_r)\right].
\]
This cancels the endpoint zero in \(Z_r(q)\).  Therefore the integrand in
\eqref{eq:generic-r-regular-integral} is nonsingular:
\begin{equation}
\frac{(\Omega-k\cos q)^2-(\Omega-c_r k)^2}{Z_r(q)}
=
\frac{-2\Omega k+k^2(\cos q+c_r)}{r+1-\cos q}.
\end{equation}
Define
\begin{equation}
\label{eq:Jr-def}
J_r
\equiv
\int_{-q_r}^{q_r}\frac{dq}{r+1-\cos q}
=
\frac{4}{\sqrt{r(r+2)}}\arccos\frac{r}{2}.
\end{equation}
Then the regular endpoint contribution is
\begin{equation}
\label{eq:generic-r-regular-final}
\Delta_{r,\rm reg}^{(--)}(\Omega,k)
=
B_{\Omega k}^{\rm reg}(r)\,\Omega k
-A_k^{\rm reg}(r)\,k^2,
\end{equation}
with
\begin{equation}
\label{eq:generic-r-regular-coeffs}
B_{\Omega k}^{\rm reg}(r)=\frac{J_r}{4\pi},
\qquad
A_k^{\rm reg}(r)=\frac{rJ_r-q_r}{4\pi}.
\end{equation}
There is no separate \(\Omega^2\) regular coefficient in this convention,
because the finite analytic term proportional to \((\Omega-c_r k)^2\) has been
absorbed into \(C_r\).  Choosing a different \(C_r\) shifts the coefficients of
\(\Omega^2\), \(\Omega k\), and \(k^2\) in the fixed ratio carried by
\((\Omega-c_r k)^2\), but it does not change
\eqref{eq:generic-r-regular-integral}--\eqref{eq:generic-r-regular-coeffs}.

Combining the singular and regular endpoint pieces in this reduced endpoint
convention gives the local large-\(\Gamma\) result
\begin{equation}
\label{eq:generic-r-full-local-endpoint}
\Delta_r^{(--)}(\Omega,k)
=
-\frac{(\Omega-c_r k)^2}{8\pi s_r}
\log\!\frac{C_r^{\rm end}}{\Gamma|\Omega-c_r k|}
+B_{\Omega k}^{\rm reg}(r)\,\Omega k
-A_k^{\rm reg}(r)\,k^2
+\cdots.
\end{equation}
The first term has been written with \(C_r^{\rm end}\) to emphasize that this
constant belongs to the reduced endpoint model with
\eqref{eq:generic-r-regular-final} kept separately.  It can be computed
analytically.  The cleanest definition is obtained on the \(k=0\) line, where
the unexpanded finite-\(\Gamma\) kernel is
\begin{equation}
\label{eq:Cr-end-kernel}
\Delta_{\rm end}(\Omega,0)
=
-\frac{\Omega^2}{8\pi}
\int_{-q_r}^{q_r}dq\,
\frac{Z_r(q)}
{Z_r(q)^2+a^2},
\qquad
a\equiv \frac{\Gamma|\Omega|}{4}.
\end{equation}
Define
\begin{equation}
\label{eq:Ia-def}
I_r(a)\equiv
\int_{-q_r}^{q_r}dq\,
\frac{Z_r(q)}{Z_r(q)^2+a^2}.
\end{equation}
As \(a\to0\), the logarithm comes only from the two endpoints.  Excise a small
distance \(\epsilon\) from both endpoints.  In the endpoint regions,
\(Z_r\simeq 2s_r p\), so the two endpoint contributions give
\begin{equation}
\label{eq:Cr-endpoint-piece}
2\int_0^\epsilon dp\,
\frac{2s_r p}{(2s_r p)^2+a^2}
=
\frac{1}{s_r}\log\frac{2s_r\epsilon}{a}+O(a^2).
\end{equation}
On the remaining interval one may set \(a=0\).  Using
\[
\frac{1}{Z_r(q)}
=
\frac{1}{2}\left[
\frac{1}{\cos q-c_r}
+\frac{1}{r+1-\cos q}
\right],
\]
the finite bulk part is
\begin{equation}
\label{eq:Cr-bulk-piece}
\int_{-q_r+\epsilon}^{q_r-\epsilon}\frac{dq}{Z_r(q)}
=
\frac{1}{s_r}\log\frac{2s_r}{\epsilon}
+\frac{J_r}{2}
+O(\epsilon),
\end{equation}
where \(J_r\) was defined in \eqref{eq:Jr-def}.  Adding
\eqref{eq:Cr-endpoint-piece} and \eqref{eq:Cr-bulk-piece} cancels the auxiliary
cutoff \(\epsilon\) and gives
\begin{equation}
\label{eq:Ia-asymptotic}
I_r(a)
=
\frac{1}{s_r}\log\frac{4s_r^2}{a}
+\frac{J_r}{2}
+o(1)
=
\frac{1}{s_r}
\log\frac{16s_r^2\exp(s_rJ_r/2)}{\Gamma|\Omega|}
+o(1).
\end{equation}
Thus, in the endpoint convention used above,
\begin{equation}
\label{eq:Cr-end-result}
C_r^{\rm end}
=
16s_r^2\exp\!\left(\frac{s_rJ_r}{2}\right).
\end{equation}
This is not a one-parameter fit to the full lattice bubble.  It is the matching
constant of the reduced endpoint kernel after the regular terms in
\eqref{eq:generic-r-regular-final} are kept separately.

The expansion \eqref{eq:generic-r-full-local-endpoint} is a small-external-
argument expansion at fixed \(\Gamma\).  Its actual control parameters are
\[
\Gamma|\Omega-c_r k|,
\qquad
\Gamma|k|,
\]
not only \(|\Omega|\) and \(|k|\).  Thus the quadratic/logarithmic coefficients
above describe the regime \(\Gamma Q\ll1\), where \(Q\) denotes a typical
external scale.  If \(\Gamma\) is taken large first at fixed nonzero
\((\Omega,k)\), the endpoint kernel crosses over to a different
\(O(\Gamma^{-2})\) form.  This noncommutativity of limits is derived in
Appendix~\ref{app:large-gamma-endpoint}.  For the full retarded-advanced
Green's function, the endpoint saturation is supplemented by smooth analytic
terms from the nonsingular components; the resulting complex kernel is given in
\eqref{eq:full-largeGamma-KRA-app}.

The same constant is obtained at nonzero \(k\).  Define
\begin{equation}
\delta\equiv \Omega-c_r k,
\qquad
D(q)\equiv \Omega-k\cos q
=\delta-k(\cos q-c_r).
\end{equation}
The finite-\(\Gamma\) endpoint kernel can be written as
\begin{equation}
\label{eq:Cr-nonzero-k-kernel}
\Delta_{\rm end}(\Omega,k)
=
-\int_{-q_r}^{q_r}\frac{dq}{2\pi}\,
\frac{Z_r(q)D(q)^2}
{4Z_r(q)^2+\Gamma^2D(q)^2/4}.
\end{equation}
Near either endpoint, \(Z_r\simeq2s_rp\) and
\(\cos q-c_r=O(p)\), so \(D(q)=\delta+O(kp)\).  The logarithmic piece is
therefore obtained by replacing \(D(q)\) by \(\delta\) inside the singular
endpoint region.  Equivalently, split
\[
D(q)^2=\delta^2+\left[D(q)^2-\delta^2\right].
\]
The first term gives the same integral \(I_r(a)\), with
\(a=\Gamma|\delta|/4\), and hence
\begin{equation}
-\frac{\delta^2}{8\pi}I_r\!\left(\frac{\Gamma|\delta|}{4}\right)
=
-\frac{\delta^2}{8\pi s_r}
\log\frac{C_r^{\rm end}}{\Gamma|\delta|}
+o(Q^2).
\end{equation}
The second term is regular, since
\begin{equation}
D(q)^2-\delta^2
=
(\cos q-c_r)\left[-2\Omega k+k^2(\cos q+c_r)\right],
\end{equation}
and the factor \(\cos q-c_r\) cancels the endpoint zero in \(Z_r(q)\).  Thus
the nonzero-\(k\) computation gives the same \(C_r^{\rm end}\) once
\(\Delta_{r,\rm reg}^{(--)}\) is kept separately.  The only exception is the
tangent direction \(\delta=\Omega-c_rk=0\): there the leading endpoint mismatch
vanishes, the \(O(kp)\) part of \(D(q)\) regulates the endpoint, and one should
use the unexpanded kernel \eqref{eq:Cr-nonzero-k-kernel} rather than extracting
a matching constant from the logarithmic formula alone.

Equivalently, if one extracts effective coefficients on a fixed small fitting
window,
\begin{align}
\Delta_r(0,k)&\simeq -\widehat A_k(r,\Gamma)\,k^2,
\\
\Delta_r(\Omega,0)&\simeq -\widehat A_\Omega(r,\Gamma)\,\Omega^2,
\\
\Delta_{r,\rm odd\text{-}odd}(\Omega,k)&\simeq
\widehat A_{\Omega k}(r,\Gamma)\,\Omega k,
\end{align}
then the logarithmic \(\Gamma\)-drifts are
\begin{align}
\frac{\partial \widehat A_k}{\partial\log\Gamma}
&=-L_k(r),
&
L_k(r)&=\frac{c_r^2}{8\pi s_r},
\label{eq:Lk-generic-r}
\\
\frac{\partial \widehat A_\Omega}{\partial\log\Gamma}
&=-L_\Omega(r),
&
L_\Omega(r)&=\frac{1}{8\pi s_r},
\label{eq:Lomega-generic-r}
\\
\frac{\partial \widehat A_{\Omega k}}{\partial\log\Gamma}
&=-L_{\Omega k}(r),
&
L_{\Omega k}(r)&=\frac{c_r}{4\pi s_r}.
\label{eq:Lomegak-generic-r}
\end{align}
The constants left after adding \(L_i(r)\log\Gamma\) are not universal: they
depend on the fitting window and on regular parts of the full Green's function.
The slopes \eqref{eq:Lk-generic-r}--\eqref{eq:Lomegak-generic-r}, however, are
fixed by the edge-bulk endpoint.

The \(r=1\) point is special because \(c_r=0\).  Then
\begin{equation}
L_k(1)=0,
\qquad
L_{\Omega k}(1)=0,
\qquad
L_\Omega(1)=\frac{1}{8\pi},
\end{equation}
and the finite endpoint-regular integrals reduce to
\begin{equation}
A_k^{\rm reg}(1)=\frac{1}{8\pi}\int_{-\pi/2}^{\pi/2}
\frac{\cos q}{2-\cos q}\,dq
=
\frac{1}{3\sqrt3}-\frac18,
\end{equation}
\begin{equation}
B_{\Omega k}^{\rm reg}(1)=\frac{1}{4\pi}\int_{-\pi/2}^{\pi/2}
\frac{dq}{2-\cos q}
=
\frac{1}{3\sqrt3}.
\end{equation}
Thus the previously used constants are recovered only at the symmetric value
\(r=1\).  For a generic value of \(r\), the leading singular kernel is instead
organized by the single endpoint combination \(\Omega-(r-1)k\).

\subsection{Finite-\texorpdfstring{$\Gamma$}{Gamma} infrared cutoffs}
\label{subsec:finite-gamma-ir-cutoffs}

There are two distinct infrared cutoffs at finite \(\Gamma\).  The first one is
the endpoint boundary layer in the reduced edge-bulk bubble.  Write
\[
\delta\equiv \Omega-c_r k.
\]
The logarithm in \eqref{eq:generic-r-singular-kernel} comes from the interval
where the distance \(p\) from an edge-bulk touching point obeys
\[
p_{\min}\ll p\ll1,
\qquad
p_{\min}\sim \frac{\Gamma|\delta|}{2s_r},
\]
up to order-one constants.  The on-shell residue approximation
\(\Sigma_R(\omega,q)\simeq Z_r(q)/(\omega-\sin q+i0)\) is itself uniform only
outside the narrower endpoint strip \(p\lesssim \Gamma^{-1}\).  Thus the clean
logarithmic matching window is
\begin{equation}
\label{eq:clean-endpoint-window}
\frac{1}{\Gamma^2}\ll |\delta|\ll \frac{1}{\Gamma}.
\end{equation}
In this window \(p_{\min}\gg\Gamma^{-1}\), and the singular term is controlled
by
\[
-\frac{\delta^2}{8\pi s_r}\log\frac{C_r}{\Gamma|\delta|}.
\]
If instead
\begin{equation}
\label{eq:ultrasmall-delta-window}
|\delta|\ll \Gamma^{-2},
\end{equation}
then the formal cutoff \(p_{\min}\) lies inside the strip where the edge pole is
already too close to the bulk branch point.  In that ultra-small-\(\delta\)
regime one should not extract the matching constant \(C_r\) from the
on-shell-residue formula.  The logarithm is cut off by the branch-point strip at
\(p\sim\Gamma^{-1}\), so the remaining \(\delta\)-dependence is analytic to this
order, with a coefficient containing \(\log\Gamma\) rather than
\(\log[1/(\Gamma|\delta|)]\).  A direct full-Green probe at \(r=1\) and
\(\Gamma=20\) indeed shows that the extracted \(\log C\) drifts once
\(\Gamma^2|\delta|=O(1)\), but this is a delicate boundary-layer effect rather
than a robust way to determine \(C_r\).

The second cutoff is the residual static bubble if the exact finite-\(\Gamma\)
local saddle normalization is replaced by its asymptotic value.  Define
\begin{equation}
\label{eq:finite-gamma-mass-def}
m_\Gamma(r)\equiv 2-\Gamma\Pi_{RA}(0,0;r,\Gamma).
\end{equation}
The exact quadratic NLSM kernel uses
\(\Pi_{RA}(\Omega,k)-\Pi_{RA}(0,0)\), so this mass is absent when the finite
\(\Gamma\) saddle cancellation is kept exactly.  If one instead substitutes
\(\Gamma\Pi_{RA}(0,0)\to2\), the diagonal kernel contains the residual local
piece
\begin{equation}
\Gamma\left[\Pi_{RA}(\Omega,k)-\frac{2}{\Gamma}\right]
=
-m_\Gamma(r)
+\Gamma\left[\Pi_{RA}(\Omega,k)-\Pi_{RA}(0,0)\right].
\end{equation}
The scaling of \(m_\Gamma\) is
\begin{equation}
\label{eq:finite-gamma-mass-scale}
m_\Gamma(r)=\frac{\mu_r}{\Gamma^2}+O(\Gamma^{-3}),
\qquad
\mu_r=O(1).
\end{equation}
This follows from the spectral-leakage identity derived in
\eqref{eq:pi00-mass-identity-app}: the defect of the monitored spectral sum
rule is supported only on the finite bulk continuum, where
\(G_R=O(\Gamma^{-1})\).  Numerically, for \(r=1\), evaluating the leakage
integral gives
\[
\begin{array}{c|ccccc}
\Gamma&20&40&80&160&320\\
\hline
\Gamma^2 m_\Gamma&1.83&2.11&2.27&2.36&2.41
\end{array}
\]
which is consistent with \eqref{eq:finite-gamma-mass-scale}.

This residual mass produces a correlation scale by comparison with the static
stiffness:
\[
A_k(r)k^2\sim m_\Gamma(r)
\qquad\Rightarrow\qquad
k_m\sim \frac{1}{\Gamma}\sqrt{\frac{\mu_r}{A_k(r)}}.
\]
Thus, if \(m_\Gamma\) is physically retained, the mass cutoff enters at
\(Q\sim\Gamma^{-1}\), whereas the ultra-small endpoint boundary layer begins
only at \(|\delta|\sim\Gamma^{-2}\).  The mass therefore cuts off the long-
wavelength physics parametrically earlier.  The regime
\(|\delta|\ll\Gamma^{-2}\) is physically relevant only for an exactly massless
kernel in which the finite-\(\Gamma\) local term has been subtracted, or tuned,
with the same accuracy as \(\Pi_{RA}(0,0)\).

\subsection{Numerical checks at representative \texorpdfstring{$r$}{r}}

We checked the logarithmic slopes using the existing direct quadrature code.
For the \(k\) and \(\Omega\) line cuts we used
\texttt{src/collect\_small\_argument\_direct\_lines.jl} with
\[
N_\omega=8192,\qquad N_q=2048,\qquad \omega_{\max}=640,
\]
on the windows
\[
5\times10^{-4}\le k\le 2.5\times10^{-3},
\qquad
5\times10^{-4}\le \Omega\le 1.5\times10^{-3}.
\]
For the mixed term we used
\texttt{src/estimate\_edge\_Aomegak.jl} at
\(\Omega=k=10^{-3}\) with the same grid.  For \(r=0.5\) and \(r=1.5\),
\[
s_r=\frac{\sqrt3}{2},
\qquad
L_k=0.0114860,
\qquad
L_\Omega=0.0459441,
\qquad
L_{\Omega k}=\mp 0.0459441,
\]
with the minus sign for \(r=0.5\) and the plus sign for \(r=1.5\).

\begin{center}
\begin{tabular}{ccccc}
\toprule
\(r\) & \(\Gamma\) &
\(\widehat A_k+L_k\log\Gamma\) &
\(\widehat A_\Omega+L_\Omega\log\Gamma\) &
\(\widehat A_{\Omega k}+L_{\Omega k}\log\Gamma\) \\
\midrule
0.5 & 40 & 0.13754 & 0.51452 & -0.14710 \\
0.5 & 80 & 0.12969 & 0.49590 & -0.14300 \\
1.5 & 40 & 0.17189 & 0.44960 & 0.56720 \\
1.5 & 80 & 0.17021 & 0.45363 & 0.56814 \\
\bottomrule
\end{tabular}
\end{center}

The shifted combinations are approximately constant between
\(\Gamma=40\) and \(80\), while the unshifted fitted coefficients drift with the
sign and magnitude predicted by \eqref{eq:Lk-generic-r}--\eqref{eq:Lomegak-generic-r}.
For example, at \(r=1.5\) the raw line-cut coefficients change as
\[
\widehat A_k: 0.12952\to0.11988,
\qquad
\widehat A_\Omega: 0.28012\to0.25230,
\qquad
\widehat A_{\Omega k}: 0.39772\to0.36681,
\]
when \(\Gamma\) is doubled from \(40\) to \(80\), consistent with a negative
\(\log\Gamma\) drift.  The \(r=0.5\) mixed coefficient has the opposite
logarithmic slope because \(L_{\Omega k}<0\), and the shifted values show the
same stabilization.  These checks support the generic-\(r\) endpoint prediction
while preserving the earlier \(r=1\) constants as a special case.

We also verified the regular coefficients
\eqref{eq:generic-r-regular-coeffs} directly in the reduced endpoint model using
\texttt{src/verify\_generic\_r\_regular\_terms.jl}.  The numerical test
integrates the finite-\(\Gamma\) endpoint kernel
\[
\Delta_{\rm end}(\Omega,k)
=
\int_{-q_r}^{q_r}\frac{dq}{2\pi}\,
\frac{-Z_r(q)(\Omega-k\cos q)^2}
{4Z_r(q)^2+\Gamma^2(\Omega-k\cos q)^2/4},
\]
subtracts the logarithmic term in \eqref{eq:generic-r-singular-kernel}, fixes
the one remaining \((\Omega-c_r k)^2\) constant from the \(\Omega\)-line, and
then extracts \(A_k^{\rm reg}\) and \(B_{\Omega k}^{\rm reg}\).  At
\(\Gamma=80\) the comparison is
\begin{center}
\begin{tabular}{ccccc}
\toprule
\(r\) &
\(A_{k,\rm num}^{\rm reg}\) &
\(A_k^{\rm reg}\) &
\(B_{\Omega k,\rm num}^{\rm reg}\) &
\(B_{\Omega k}^{\rm reg}\) \\
\midrule
0.5 & 0.02197 & 0.02097 & 0.37215 & 0.37527 \\
1.0 & 0.06734 & 0.06745 & 0.19028 & 0.19245 \\
1.5 & 0.06570 & 0.06727 & 0.09689 & 0.10040 \\
\bottomrule
\end{tabular}
\end{center}
The agreement is at the few-percent level on the finite fitting window used for
the check.  The full Green's-function quadrature verifies the logarithmic slopes
cleanly; its absolute regular constants also contain nonsingular lattice
contributions, so the reduced endpoint table is the sharper test of the
analytic regular terms in \eqref{eq:generic-r-regular-coeffs}.

The analytical constant \(C_r^{\rm end}\) in \eqref{eq:Cr-end-result} is only
the reduced endpoint matching constant.  It is not, by itself, a verification of
the full lattice Green's function.  The full Green's function contains
nonsingular regular terms that are also \(O(\Omega^2)\) on the \(k=0\) line, so
the constant multiplying the logarithm should be extracted from the numerical
full bubble after the logarithmic slope has been fixed.

For this purpose it is better to work with the logarithm of the constant.  On
the \(k=0\) line define
\begin{equation}
\label{eq:lnCr-full-estimator}
\ell_r^{\rm full}(\Omega;\Gamma)
\equiv
\log C_{r,\rm est}^{\rm full}
=
\log(\Gamma|\Omega|)
-
\frac{
\Delta_{\rm full}(\Omega,0)}
{L_\Omega(r)\Omega^2},
\qquad
L_\Omega(r)=\frac{1}{8\pi s_r},
\end{equation}
where
\[
\Delta_{\rm full}(\Omega,0)
=
\Gamma\,\Re\!\left[
\frac{\Pi_{RA}(\Omega,0)+\Pi_{RA}(-\Omega,0)}{2}
-\Pi_{RA}(0,0)\right]
\]
is computed from the full orbital Green's function.  In this convention all
regular \(O(\Omega^2)\) terms of the full bubble are absorbed into
\(\ell_r^{\rm full}\).  Thus \(\ell_r^{\rm full}\) should be compared with
\(\log C_r^{\rm end}\), not with \(C_r^{\rm end}\) itself.

This extraction is numerically more delicate than the logarithmic-slope check,
because it divides a small \(O(\Omega^2)\) difference by \(\Omega^2\).  In
particular, the frequency mesh must resolve the external frequency itself.  A
coarse dense-region spacing \(d\omega=1/(20\Gamma)\) gives an artificially low
\(\ell_r^{\rm full}\) near \(r=1\).  We therefore use a refined spacing
\(d\omega=1/(100\Gamma)\) for the quoted full-Green values.

We performed this extraction with \texttt{src/estimate\_full\_Cr.jl}.  The
script evaluates the full orbital trace and, in the same run, the \(\sigma_x\)
\(--\) component.  The raw complex bubbles and grids are saved in
\texttt{data/full\_green\_lnCr\_refined.jld2}.  The table uses
\(\Gamma=80\), \(\Omega=1.5\times10^{-3}\), and hence
\(\Gamma\Omega=0.12\):
\begin{center}
\begin{tabular}{ccccc}
\toprule
\(r\) &
\(\log C_r^{\rm end}\) &
\(\ell_r^{\rm full}(1.5\times10^{-3})\) &
\(\ell_r^{--}(1.5\times10^{-3})\) &
\(\ell_r^{\rm full}-\log C_r^{\rm end}\) \\
\midrule
0.5 & 4.52692 & 4.52236 & 4.51515 & \(-0.00456\) \\
1.0 & 3.98179 & 3.95438 & 3.94670 & \(-0.02740\) \\
1.5 & 3.03124 & 3.00945 & 3.00178 & \(-0.02179\) \\
\bottomrule
\end{tabular}
\end{center}
The full result is indeed dominated by the \(\sigma_x\) \({--}\) bubble: the
middle columns differ only at the \(10^{-2}\) level in \(\log C\).  The
residual difference from \(\log C_r^{\rm end}\) is much smaller than what one
finds on the coarse frequency mesh, and is comparable to remaining
finite-\(\Gamma\Omega\) and finite-mesh uncertainties.

For reference, at the most sensitive point \(r=1\) the frequency convergence is
\begin{center}
\begin{tabular}{cccc}
\toprule
\(d\omega\) & \(N_\omega\) &
\(\ell_{1}^{\rm full}(1.5\times10^{-3})\) &
\(\ell_{1}^{--}(1.5\times10^{-3})\) \\
\midrule
\(1/(20\Gamma)\)  & 8952  & 3.86625 & 3.85870 \\
\(1/(50\Gamma)\)  & 18552 & 3.94260 & 3.93495 \\
\(1/(100\Gamma)\) & 34552 & 3.95438 & 3.94670 \\
\(1/(200\Gamma)\) & 66552 & 3.95900 & 3.95134 \\
\bottomrule
\end{tabular}
\end{center}
This convergence data is saved in
\texttt{data/full\_green\_lnCr\_r1\_frequency\_convergence.jld2}.  The main
lesson is that the full integral is dominated by the \({--}\) bubble, but the
matching constant is not a robust quantity unless the regular \(O(\Omega^2)\)
terms and the numerical frequency resolution are controlled.  The logarithmic
slopes are much more reliable than the absolute constant.

\section{Trivial-\texorpdfstring{$r>2$}{r>2} Bubble}
\label{sec:trivial-r-bubble}

For \(r>2\), the boundary is in the trivial phase.  The condition
\(|r-\cos q|<1\) has no solution, so the edge-pole interval
\(\mathcal I_r\) is empty.  The endpoint mechanism responsible for
\(\Omega^2\log|\Omega|\), \(k^2\log\Gamma\), and
\(\Omega k\log\Gamma\) is therefore absent.  The full retarded-advanced bubble
is analytic at small \((\Omega,k)\).

It is convenient to fit the unsymmetrized full bubble as
\begin{equation}
\label{eq:trivial-fit-def}
\Gamma\Pi_{RA}(\Omega,k)
=
C_0
+L_\Omega\Omega
+L_k k
+Q_{\Omega\Omega}\Omega^2
+Q_{\Omega k}\Omega k
+Q_{kk}k^2
+\cdots .
\end{equation}
The large-\(\Gamma\) expansion is controlled by the broadening-dominated Green's
function.  Since the trivial bulk is gapped, the self-energy is analytic near
zero frequency and obeys \(\Sigma_R(z,q)=1/z+O(z^{-2})\) at large \(|z|\).
For \(\omega=\Gamma x/2\),
\begin{equation}
G_R\!\left(\frac{\Gamma x}{2}+\Omega,q+k\right)
=
\frac{2}{\Gamma}\frac{\mathbf{1}}{x+i+2\Omega/\Gamma}
+
\frac{4}{\Gamma^2}\frac{h_r(q+k)}{(x+i)^2}
+O(\Gamma^{-3}),
\end{equation}
where
\[
h_r(q)=\sin q\,\sigma_x+(r-\cos q)\,\sigma_z.
\]
The scalar part gives
\begin{equation}
\label{eq:trivial-scalar-expansion}
\Gamma\Pi_{RA}^{(0)}(\Omega,k)
=
\frac{2}{1-i\Omega/\Gamma}
=
2+\frac{2i\Omega}{\Gamma}
-\frac{2\Omega^2}{\Gamma^2}
+O(\Gamma^{-3}Q^3).
\end{equation}
The first momentum-dependent term comes from \(h_r(q+k)h_r(q)\).  Subtracting
the \(k=0\) value,
\begin{align}
\Gamma\left[\Pi_{RA}(0,k)-\Pi_{RA}(0,0)\right]_{h^2}
&=
\frac{2}{\Gamma^2}
\int_{-\pi}^{\pi}\frac{dq}{2\pi}\,
\tr\!\left[h_r(q+k)h_r(q)-h_r(q)^2\right]
\notag\\
&=
\frac{4(\cos k-1)}{\Gamma^2}
=
-\frac{2k^2}{\Gamma^2}+O(k^4/\Gamma^2).
\label{eq:trivial-k2-derivation}
\end{align}
The trace of \(h_r\) vanishes and the scalar denominator is independent of
\(q\), so there is no linear \(k\) term and no \(\Omega k\) term at this order.

The constant term is most cleanly obtained from the spectral-leakage identity
\eqref{eq:pi00-mass-identity-app}.  In the trivial phase the continuum exhausts
the self-energy spectral weight:
\[
\int\frac{d\omega}{2\pi}\,
\left[-2\,\mathrm{Im}\,\Sigma_R(\omega,q)\right]=1
\]
for each \(q\).  Since
\(\tr[G_RP_-G_A]=4/\Gamma^2+O(\Gamma^{-3})\) on the finite bulk continuum,
\begin{equation}
\label{eq:trivial-C0-result}
C_0(r,\Gamma)
=
\Gamma\Pi_{RA}(0,0)
=
2-\frac{4}{\Gamma^2}+O(\Gamma^{-3}),
\qquad r>2.
\end{equation}
Combining these pieces gives the leading trivial-phase coefficients
\begin{equation}
\label{eq:trivial-coeff-summary}
\begin{gathered}
C_0=2-\frac{4}{\Gamma^2}+O(\Gamma^{-3}),
\qquad
L_\Omega=\frac{2i}{\Gamma}+O(\Gamma^{-3}),
\qquad
L_k=O(\Gamma^{-3}),
\\
Q_{\Omega\Omega}=-\frac{2}{\Gamma^2}+O(\Gamma^{-3}),
\qquad
Q_{kk}=-\frac{2}{\Gamma^2}+O(\Gamma^{-3}),
\qquad
Q_{\Omega k}=O(\Gamma^{-3}).
\end{gathered}
\end{equation}
The leading coefficients are independent of \(r\) as long as \(r>2\).  The
\(r\)-dependence first enters through subleading corrections to the analytic
large-\(\Gamma\) expansion.

We checked these predictions at \(r=3\) by direct full-Green quadrature using
\texttt{src/analyze\_trivial\_phase\_bubble.jl}.  The raw bubbles and fit
outputs are saved in \texttt{data/trivial\_r3\_bubble\_scan.jld2}.  The table
shows the scaled coefficients; the first four columns should approach
\(4,2,-2,-2\), while the last two columns remain \(O(1)\) because the
corresponding coefficients start only at \(O(\Gamma^{-3})\).
\begin{center}
\begin{tabular}{ccccccc}
\toprule
\(\Gamma\) &
\(\Gamma^2m_{\rm leak}\) &
\(\Gamma\,\Im L_\Omega\) &
\(\Gamma^2\Re Q_{\Omega\Omega}\) &
\(\Gamma^2\Re Q_{kk}\) &
\(\Gamma^3\Im L_k\) &
\(\Gamma^3\Re Q_{\Omega k}\) \\
\midrule
20  & 2.95290 & 2.00183 & -3.12290 & -2.63531 & -3.97411 & 15.15975 \\
40  & 3.50134 & 2.00078 & -2.65491 & -2.47261 & -5.46633 & 21.94952 \\
80  & 3.77522 & 2.00011 & -2.26667 & -2.28032 & -6.22272 & 24.90956 \\
160 & 3.89637 & 1.99988 & -2.03910 & -2.15607 & -6.56212 & 25.80893 \\
\bottomrule
\end{tabular}
\end{center}
Here \(m_{\rm leak}\) denotes the positive residual
\(2-\Gamma\Pi_{RA}(0,0)\) evaluated through
\eqref{eq:pi00-mass-identity-app}, which is more stable than subtracting the
tail-corrected constant at the largest \(\Gamma\).  The numerical trend confirms
that the trivial phase has no separate \(\Gamma Q\ll1\) endpoint regime.  After
the constant piece is subtracted, the small-\((\Omega,k)\) kernel decays
algebraically with \(\Gamma\), with the leading terms shown in
\eqref{eq:trivial-coeff-summary}.

\section{One-loop RG of the replica sigma model}

We now use the replica-sector kernel to identify the sigma-model stiffness and
to derive its one-loop RG flow.  Since the full KK bubble is smooth and strongly
suppressed at large \(\Gamma\), we keep only the retarded-advanced part in the
replica sector.  From \eqref{eq:Y-sector-general}, \eqref{eq:Y-kernel-general},
and \eqref{eq:Delta-summary-main}, the Euclidean quadratic action of the
traceless replica mode is
\begin{equation}\label{eq:SE-Y-quadratic}
S_E^{(2)}[Y]
=
\frac{\Gamma}{4}\int_Q
\left[
A_k\,k^2
-A_{\Omega k}\,\Omega k
+\frac{\Omega^2}{8\pi}\log\!\frac{C_\Omega}{\Gamma|\Omega|}
+\cdots
\right]
\Tr_R\!\left[Y(-Q)Y(Q)\right].
\end{equation}
Matching the spatial-gradient term to the standard normalization
\(\frac{\lambda}{8}\int_Q k^2 \Tr_R[Y(-Q)Y(Q)]\), we identify
\begin{equation}\label{eq:lambda-def}
\frac{\lambda}{8}\equiv \frac{\Gamma A_k}{4},
\qquad
\lambda=2\Gamma A_k.
\end{equation}
Using the large-\(\Gamma\) value of \(A_k\),
\begin{equation}\label{eq:lambda-micro}
\lambda
\xrightarrow[\Gamma\to\infty]{}
\Gamma\left(\frac{2}{3\sqrt{3}}-\frac{1}{4}\right)
=
\Gamma\,\frac{8-3\sqrt{3}}{12\sqrt{3}}.
\end{equation}

It is also convenient to define
\begin{equation}\label{eq:v-rho-def}
v\equiv \frac{A_{\Omega k}}{A_k},
\qquad
\kappa\equiv \frac{1}{8\pi A_k},
\qquad
\rho(\Omega)\equiv \kappa\log\!\frac{C_\Omega}{\Gamma|\Omega|}.
\end{equation}
Then \eqref{eq:SE-Y-quadratic} becomes
\begin{equation}\label{eq:SE-Y-kernel}
S_E^{(2)}[Y]
=
\frac{\lambda}{8}\int_Q
\left[
k^2-v\,\Omega k+\rho(\Omega)\,\Omega^2+\cdots
\right]
\Tr_R\!\left[Y(-Q)Y(Q)\right].
\end{equation}
The mixed term can be removed by the shear
\begin{equation}
\tilde k = k-\frac{v}{2}\Omega,
\end{equation}
under which
\begin{equation}\label{eq:kernel-sheared}
k^2-v\,\Omega k+\rho(\Omega)\Omega^2
=
\tilde k^{\,2}
+\tilde\rho(\Omega)\,\Omega^2,
\qquad
\tilde\rho(\Omega)\equiv \rho(\Omega)-\frac{v^2}{4}.
\end{equation}
More generally, the omitted analytic \(\Omega^2\) terms in
\eqref{eq:SE-Y-quadratic} can be absorbed into a regular constant \(r\), so the
time kernel takes the form
\begin{equation}\label{eq:rho-decomp}
\tilde\rho(\Omega)=r+\kappa\log\!\frac{\Omega_0}{|\Omega|},
\qquad
\Omega_0\equiv \frac{C_\Omega}{\Gamma}.
\end{equation}
At asymptotically small \(|\Omega|\), the logarithm wins over the regular part,
so \(\tilde\rho(\Omega)\) remains positive.  We therefore perform the RG in the
sheared frame and drop the tilde on \(k\) from now on.

The nonlinear completion of the replica sector is obtained by introducing
\begin{equation}\label{eq:X-field-def}
X(\tau,x)\in SU(R),
\qquad
J_\mu\equiv X^{-1}\partial_\mu X,
\end{equation}
so that \(J_\mu=i\partial_\mu Y+O(Y^2)\).  The replica sigma model is then
\begin{equation}\label{eq:SE-X}
S_E[X]
=
\frac{\lambda}{8}
\left[
\int d\tau\,dx\,\Tr_R\!\left(J_x^2\right)
+\int_Q \rho(\Omega)\,
\Tr_R\!\left(J_\tau(-Q)J_\tau(Q)\right)
\right].
\end{equation}
The logarithm therefore remains attached to the time-derivative current
bilinear; it is not replaced by a local \((\partial_\tau X)^2\) term.

To carry out Wilsonian RG, split
\begin{equation}\label{eq:X-split}
X=X_s X_f,
\qquad
X_f=e^{iY_f},
\end{equation}
where \(Y_f\) contains only shell modes.  Writing
\begin{equation}
J_\mu^{(s)}\equiv X_s^{-1}\partial_\mu X_s,
\end{equation}
we have
\begin{equation}
J_\mu
=
X_f^{-1}J_\mu^{(s)}X_f+X_f^{-1}\partial_\mu X_f.
\end{equation}
Expanding \(X_f=e^{iY_f}\) gives
\begin{equation}\label{eq:J-expand-fast}
J_\mu
=
J_\mu^{(s)}
+i\nabla_\mu Y_f
-\frac{1}{2}[Y_f,\nabla_\mu Y_f]
-\frac{1}{2}[Y_f,[Y_f,J_\mu^{(s)}]]
+O(Y_f^3),
\end{equation}
with the background covariant derivative
\begin{equation}
\nabla_\mu Y_f \equiv \partial_\mu Y_f + i[J_\mu^{(s)},Y_f].
\end{equation}
Substituting \eqref{eq:J-expand-fast} into \eqref{eq:SE-X}, the linear term in
\(Y_f\) vanishes after integrating by parts.  The quadratic fast action is
\begin{equation}\label{eq:S-fast}
S_f[Y_f]
=
\frac{\lambda}{8}\int_Q
\mathcal K(Q)\,
\Tr_R\!\left[Y_f(-Q)Y_f(Q)\right],
\qquad
\mathcal K(Q)=k^2+\rho(\Omega)\Omega^2.
\end{equation}
The logarithm must therefore be kept inside the fast-mode propagator throughout
the shell integration.

Choose a basis \(T^a\) of \(su(R)\) with
\begin{equation}
\Tr_R(T^a T^b)=\frac{1}{2}\delta^{ab},
\qquad
Y_f = Y_f^a T^a .
\end{equation}
Then
\begin{equation}\label{eq:Yf-prop}
\left\langle
Y_f^a(Q)Y_f^b(-Q)
\right\rangle_>
=
\frac{8}{\lambda}\,\delta^{ab}\,\frac{1}{\mathcal K(Q)},
\end{equation}
because
\begin{equation}
\Tr_R\!\left[Y_f(-Q)Y_f(Q)\right]
=
\frac{1}{2}Y_f^a(-Q)Y_f^a(Q).
\end{equation}
This factor of \(1/2\) from the trace normalization is the source of the
missing factor of two in the naive estimate.
and the shell average at coincident points is
\begin{equation}\label{eq:I-shell}
I_>\equiv
\int_>\frac{d\Omega\,dk}{(2\pi)^2}\,\frac{1}{\mathcal K(Q)}.
\end{equation}
For an infinitesimal Wilson shell defined by
\(\Lambda e^{-d\ell}<\sqrt{\mathcal K(Q)}<\Lambda\), the logarithm varies only
slowly across the shell.  Freezing it at the running shell value
\begin{equation}
\rho_\ell \equiv r+\kappa \ell,
\end{equation}
with \(|\Omega| \sim \Omega_0 e^{-\ell}\), we obtain
\begin{equation}\label{eq:I-shell-dl}
I_>
\simeq
\int_>\frac{d\Omega\,dk}{(2\pi)^2}\,\frac{1}{k^2+\rho_\ell \Omega^2}
=
\frac{d\ell}{2\pi\sqrt{\rho_\ell}}
+O(d\ell^2).
\end{equation}
So the logarithmic time kernel does affect the RG: it reduces the one-loop shell
weight by the factor \(1/\sqrt{\rho_\ell}\), which slowly decreases with scale.

The terms quadratic in both \(J_\mu^{(s)}\) and \(Y_f\) renormalize the slow
action.  Using
\begin{equation}
\sum_a [T^a,[T^a,M]] = R\,M
\qquad
\text{for } M\in su(R),
\end{equation}
one finds
\begin{equation}\label{eq:Seff-renorm}
S_E^{\rm eff}[X_s]
=
\frac{\lambda-2RI_>}{8}
\left[
\int d\tau\,dx\,\Tr_R\!\left((J_x^{(s)})^2\right)
+\int_Q \rho(\Omega)\,
\Tr_R\!\left(J_\tau^{(s)}(-Q)J_\tau^{(s)}(Q)\right)
\right].
\end{equation}
Thus the spatial-gradient term and the logarithmic time-derivative term are
renormalized by the same loop factor.  However, this does \emph{not} mean that
the logarithmic kernel is inert.  It affects the RG in two ways.  First, it
enters the shell integral through \eqref{eq:I-shell-dl}.  Second, under the
rescaling \(\Omega\to e^{d\ell}\Omega\),
\begin{equation}\label{eq:rho-rescale}
\rho\!\left(\Omega e^{-d\ell}\right)
=
\rho(\Omega)+\kappa\,d\ell,
\end{equation}
so the regular \(\Omega^2\) coefficient shifts at tree level.  Equivalently, if
\(\rho(\Omega)=r+\kappa\log(\Omega_0/|\Omega|)\), then
\begin{equation}\label{eq:r-kappa-flow}
\frac{d\kappa}{d\ell}=0+O(\lambda^2),
\qquad
\frac{dr}{d\ell}=\kappa+O(\lambda).
\end{equation}
At this order the loop correction does not generate a separate beta function for
\(\kappa\); it only renormalizes the overall stiffness.

Reading off the renormalized coupling from \eqref{eq:Seff-renorm},
\begin{equation}
\lambda(\ell+d\ell)
=
\lambda(\ell)-2R\,I_>,
\end{equation}
and using \eqref{eq:I-shell-dl}, we obtain
\begin{equation}\label{eq:beta-lambda}
\frac{d\lambda}{d\ell}
=
-\frac{R}{\pi\sqrt{\rho_\ell}}
+O(\lambda^{-1}).
\end{equation}
Thus the replica sector has the same one-loop algebraic structure as the
principal chiral model, but with a scale-dependent shell weight set by the
logarithmic time kernel.  In the replica continuation relevant here,
\(R\to 1\), the one-loop flow is
\begin{equation}
\beta_\lambda
=
-\frac{1}{\pi\sqrt{\rho_\ell}}
+O(\lambda^{-1}),
\end{equation}
while \(\rho_\ell=r+\kappa\ell\) grows logarithmically with the RG scale.
Consequently the logarithmic time kernel suppresses the growth of \(\lambda\) at
long distances through the factor \(1/\sqrt{\rho_\ell}\).  In the present
convention, \(\lambda\) is a true stiffness, so the negative sign of
\eqref{eq:beta-lambda} means that the stiffness decreases under coarse graining.

The one-loop flow can be integrated explicitly.  For general \(R\),
\begin{equation}\label{eq:lambda-flow-integrated}
\lambda(\ell)
=
\lambda_0
-\frac{2R}{\pi\kappa}
\left(
\sqrt{r+\kappa\ell}-\sqrt{r}
\right),
\qquad
\kappa>0.
\end{equation}
In the limit \(\kappa\to 0\), this reduces smoothly to the isotropic result
\begin{equation}\label{eq:lambda-flow-nolog}
\lambda(\ell)
=
\lambda_0-\frac{R}{\pi\sqrt{r}}\,\ell.
\end{equation}
For the standard isotropic normalization \(r=1\), one recovers
\begin{equation}
\frac{d\lambda}{d\ell}=-\frac{R}{\pi},
\qquad
\lambda(\ell)=\lambda_0-\frac{R}{\pi}\ell.
\end{equation}

The scale \(\ell_*\) at which the stiffness is driven to zero,
\(\lambda(\ell_*)=0\), estimates the onset of the disordered regime.  Without
the logarithmic term,
\begin{equation}\label{eq:ellstar-nolog}
\ell_*^{(0)}=\frac{\pi\sqrt{r}}{R}\,\lambda_0
\qquad
\left(\ell_*^{(0)}=\frac{\pi}{R}\lambda_0\ \text{for } r=1\right).
\end{equation}
With the logarithmic kernel present,
\begin{equation}\label{eq:ellstar-log}
\ell_*=
\frac{\pi\sqrt{r}}{R}\,\lambda_0
+\frac{\pi^2\kappa}{4R^2}\,\lambda_0^2.
\end{equation}
Thus the logarithmic term slows the flow toward disorder.  Instead of a constant
velocity \(d\lambda/d\ell=-R/\pi\), the magnitude of the beta function decays as
\begin{equation}
\left|\beta_\lambda\right|
\sim
\frac{R}{\pi\sqrt{\kappa\ell}}
\qquad
(\ell\gg r/\kappa),
\end{equation}
so the stiffness decreases only as \(\lambda(\ell)\sim \lambda_0-\text{const}
\times\sqrt{\ell}\).  Equivalently, the disordering scale is parametrically
larger:
\begin{equation}
\ell_*-\ell_*^{(0)}
=
\frac{\pi^2\kappa}{4R^2}\,\lambda_0^2.
\end{equation}
If one interprets \(\xi\sim e^{\ell_*}\) as the correlation length, then the
correlation length is
\begin{equation}\label{eq:xi-log}
\xi
\sim
\xi_0\,e^{\ell_*}
=
\xi_0
\exp\!\left[
\frac{\pi\sqrt{r}}{R}\lambda_0
+\frac{\pi^2\kappa}{4R^2}\lambda_0^2
\right],
\end{equation}
to be compared with the isotropic result without the logarithmic kernel,
\begin{equation}\label{eq:xi-nolog}
\xi^{(0)}
\sim
\xi_0
\exp\!\left[
\frac{\pi\sqrt{r}}{R}\lambda_0
\right].
\end{equation}
Thus
\begin{equation}\label{eq:xi-ratio}
\frac{\xi}{\xi^{(0)}}
=
\exp\!\left[
\frac{\pi^2\kappa}{4R^2}\lambda_0^2
\right].
\end{equation}
The logarithmic \(\Omega^2\) term therefore produces a parametrically larger
correlation length.  For moderate bare stiffness, where
\(\kappa\lambda_0\ll 4R\sqrt{r}/\pi\), the extra term in the exponent is a small
correction.  For sufficiently large \(\lambda_0\), however, the quadratic term
dominates, and the correlation length crosses over from the usual
\(\exp(\text{const}\times \lambda_0)\) behavior to the stronger form
\(\exp(\text{const}\times \lambda_0^2)\).

Using the microscopic large-\(\Gamma\) expressions
\begin{equation}
\lambda_0 \simeq 2\Gamma A_k,
\qquad
\kappa=\frac{1}{8\pi A_k},
\end{equation}
and taking \(R\to 1\), one obtains
\begin{equation}\label{eq:xi-log-micro}
\ell_*
\simeq
2\pi A_k\sqrt{r}\,\Gamma
+\frac{\pi A_k}{8}\Gamma^2,
\end{equation}
so that at large \(\Gamma\)
\begin{equation}
\xi
\sim
\xi_0\,
\exp\!\left[
2\pi A_k\sqrt{r}\,\Gamma
+\frac{\pi A_k}{8}\Gamma^2
\right].
\end{equation}
At the one-loop level, this is the main physical consequence of the logarithmic
time kernel: it does not stop the flow toward disorder, but it slows that flow
enough to promote the correlation length from the usual exponential dependence
to a super-exponential one in the microscopic stiffness.

\section{Generic-\texorpdfstring{$r$}{r} One-loop RG in the Ballistic Frame}
\label{sec:generic-r-rg}

The RG discussion above was written in the symmetric case \(r=1\), where the
nonanalytic piece sits on \(\Omega^2\).  For a generic topological value
\(0<r<2\), Section~\ref{sec:generic-r-bubbles} showed that the singular bubble
depends instead on the endpoint combination
\begin{equation}
\delta_r\equiv \Omega-c_r k,
\qquad
c_r=r-1,
\qquad
s_r=\sqrt{r(2-r)}.
\end{equation}
Therefore the natural low-energy frame is not the laboratory frame
\((\Omega,k)\), but the ballistic frame aligned with the edge velocity \(c_r\).

To avoid clashing with the microscopic mass parameter \(r\), in this section we
denote the regular analytic coefficient of the time kernel by \(\rho_{0,r}\)
rather than by \(r\).  The quadratic replica action consistent with
\eqref{eq:generic-r-full-local-endpoint} can be written as
\begin{equation}\label{eq:SE-Y-generic-r-lab}
S_{E,r}^{(2)}[Y]
=
\frac{\Gamma}{4}\int_Q
\left[
A_{k,r}\,k^2
-A_{\Omega k,r}\,\Omega k
+A_{\Omega,r}\,\Omega^2
+\frac{(\Omega-c_r k)^2}{8\pi s_r}
\log\!\frac{\Omega_{0,r}}{|\Omega-c_r k|}
+\cdots
\right]
\Tr_R\!\left[Y(-Q)Y(Q)\right].
\end{equation}
Here \(\Omega_{0,r}\sim C_r/\Gamma\) is the matching scale, while
\(A_{k,r}\), \(A_{\Omega k,r}\), and \(A_{\Omega,r}\) are finite analytic
coefficients.  They depend on the convention used to separate the regular and
logarithmic pieces.  In the reduced endpoint convention of
\eqref{eq:generic-r-regular-coeffs} and \eqref{eq:Cr-end-result},
\begin{equation}\label{eq:generic-r-endpoint-bare}
A_{\Omega,r}^{\rm end}=0,
\qquad
A_{\Omega k,r}^{\rm end}=B_{\Omega k}^{\rm reg}(r)=\frac{J_r}{4\pi},
\qquad
A_{k,r}^{\rm end}=A_k^{\rm reg}(r)=\frac{rJ_r-q_r}{4\pi}.
\end{equation}
The full lattice Green's function shifts these coefficients by additional
regular \(O(1)\) terms, but it does not change the universal logarithmic
coefficient \(1/(8\pi s_r)\).

The singular direction is made explicit by the unit-Jacobian change of
variables
\begin{equation}\label{eq:generic-r-ballistic-vars}
\varpi\equiv \Omega-c_r k,
\qquad
p\equiv k.
\end{equation}
In real space this is the comoving coordinate system
\begin{equation}
\tau'=\tau,
\qquad
x'=x-c_r\tau,
\end{equation}
so \(\varpi\) is conjugate to \(\tau'\), and
\begin{equation}
\partial_{\tau'}=\partial_\tau+c_r\partial_x,
\qquad
\partial_{x'}=\partial_x.
\end{equation}
Thus the logarithmic kernel is attached to the convective derivative along the
edge trajectory rather than to \(\partial_\tau\) alone.

Substituting \(\Omega=\varpi+c_r p\) into \eqref{eq:SE-Y-generic-r-lab} gives
\begin{equation}\label{eq:SE-Y-generic-r-ballistic}
S_{E,r}^{(2)}[Y]
=
\frac{\Gamma}{4}\int_Q
\left[
\bar A_{k,r}\,p^2
-\bar A_{\varpi p,r}\,\varpi p
+\left(
\bar A_{\varpi,r}
+\frac{1}{8\pi s_r}\log\!\frac{\Omega_{0,r}}{|\varpi|}
\right)\varpi^2
+\cdots
\right]
\Tr_R\!\left[Y(-Q)Y(Q)\right],
\end{equation}
with
\begin{align}
\bar A_{k,r}
&=
A_{k,r}-c_rA_{\Omega k,r}+c_r^2A_{\Omega,r},
\label{eq:generic-r-Abar-k}
\\
\bar A_{\varpi p,r}
&=
A_{\Omega k,r}-2c_rA_{\Omega,r},
\label{eq:generic-r-Abar-mix}
\\
\bar A_{\varpi,r}
&=
A_{\Omega,r}.
\label{eq:generic-r-Abar-varpi}
\end{align}
In the reduced endpoint convention \eqref{eq:generic-r-endpoint-bare}, these
become
\begin{equation}\label{eq:generic-r-endpoint-ballistic-coeffs}
\bar A_{k,r}^{\rm end}
=
A_k^{\rm reg}(r)-c_r B_{\Omega k}^{\rm reg}(r)
=
\frac{J_r-q_r}{4\pi},
\qquad
\bar A_{\varpi p,r}^{\rm end}
=
\frac{J_r}{4\pi},
\qquad
\bar A_{\varpi,r}^{\rm end}=0.
\end{equation}
Since \(J_r-q_r>0\) for \(0<r<2\), the ballistic spatial stiffness is positive
throughout the topological phase.

Matching the \(p^2\) term to the standard normalization gives
\begin{equation}\label{eq:generic-r-lambda-def}
\frac{\lambda_r}{8}\equiv \frac{\Gamma\bar A_{k,r}}{4},
\qquad
\lambda_r=2\Gamma\bar A_{k,r}.
\end{equation}
It is then natural to define
\begin{equation}\label{eq:generic-r-rho-def}
v_r\equiv \frac{\bar A_{\varpi p,r}}{\bar A_{k,r}},
\qquad
\kappa_r\equiv \frac{1}{8\pi s_r\bar A_{k,r}},
\qquad
\rho_r(\varpi)\equiv
\rho_{0,r}
+\kappa_r\log\!\frac{\Omega_{0,r}}{|\varpi|},
\qquad
\rho_{0,r}\equiv \frac{\bar A_{\varpi,r}}{\bar A_{k,r}}.
\end{equation}
Then \eqref{eq:SE-Y-generic-r-ballistic} becomes
\begin{equation}\label{eq:SE-Y-generic-r-kernel}
S_{E,r}^{(2)}[Y]
=
\frac{\lambda_r}{8}\int_Q
\left[
p^2-v_r\varpi p+\rho_r(\varpi)\varpi^2+\cdots
\right]
\Tr_R\!\left[Y(-Q)Y(Q)\right].
\end{equation}
In the reduced endpoint convention,
\begin{equation}\label{eq:generic-r-endpoint-rg-couplings}
\lambda_r^{\rm end}
=
\Gamma\,\frac{J_r-q_r}{2\pi},
\qquad
v_r^{\rm end}
=
\frac{J_r}{J_r-q_r},
\qquad
\kappa_r^{\rm end}
=
\frac{1}{2s_r(J_r-q_r)}.
\end{equation}
At \(r=1\), where \(c_r=0\), \(q_r=\pi/2\), and \(J_r=4\pi/(3\sqrt3)\), these
reduce to the \(r=1\) couplings used above.

As in \eqref{eq:kernel-sheared}, the remaining analytic mixed term is removed by
the shear
\begin{equation}
\tilde p=p-\frac{v_r}{2}\varpi,
\end{equation}
under which
\begin{equation}\label{eq:generic-r-kernel-sheared}
p^2-v_r\varpi p+\rho_r(\varpi)\varpi^2
=
\tilde p^{\,2}
+\tilde\rho_r(\varpi)\varpi^2,
\qquad
\tilde\rho_r(\varpi)
\equiv
\tilde\rho_{0,r}
+\kappa_r\log\!\frac{\Omega_{0,r}}{|\varpi|},
\end{equation}
with
\begin{equation}
\tilde\rho_{0,r}\equiv \rho_{0,r}-\frac{v_r^2}{4}.
\end{equation}
The omitted analytic \(O(\varpi^2)\) terms from the full Green's function only
shift \(\tilde\rho_{0,r}\) by a finite amount.  The logarithmic coefficient
\(\kappa_r\) is fixed by the endpoint data.  In the reduced endpoint
convention, \(\rho_{0,r}^{\rm end}=0\), so
\begin{equation}
\tilde\rho_{0,r}^{\rm end}
=
-\frac{(v_r^{\rm end})^2}{4}
=
-\frac{J_r^2}{4(J_r-q_r)^2}.
\end{equation}
This negative \(O(1)\) offset is convention-dependent: the full lattice Green's
function adds further analytic \(\varpi^2\) terms.  The RG below is therefore
understood to start at scales where the logarithm has already made
\(\tilde\rho_r(\varpi)\) positive.

The nonlinear completion is therefore the same principal-chiral replica sigma
model as before, but written in the comoving coordinates
\((\tau',x')\):
\begin{equation}\label{eq:SE-X-generic-r}
S_{E,r}[X]
=
\frac{\lambda_r}{8}
\left[
\int d\tau'\,dx'\,\Tr_R\!\left(J_{x'}^2\right)
+\int_Q \tilde\rho_r(\varpi)\,
\Tr_R\!\left(J_{\tau'}(-Q)J_{\tau'}(Q)\right)
\right],
\end{equation}
where \(J_{\mu'}\equiv X^{-1}\partial_{\mu'}X\).  In other words, the generic
\(0<r<2\) theory is not isotropic in the laboratory variables, but it becomes
the same logarithmic sigma model after passing to the ballistic frame selected
by the edge velocity \(c_r\).

Because the derivation of the fast-mode expansion depends only on the quadratic
kernel, the one-loop calculation goes through unchanged with
\((\Omega,k)\) replaced by \((\varpi,\tilde p)\).  The fast propagator is
\begin{equation}
\left\langle
Y_f^a(Q)Y_f^b(-Q)
\right\rangle_>
=
\frac{8}{\lambda_r}\,\delta^{ab}\,
\frac{1}{\tilde p^{\,2}+\tilde\rho_r(\varpi)\varpi^2},
\end{equation}
and the infinitesimal shell integral is
\begin{equation}\label{eq:generic-r-shell}
I_>^{(r)}
\equiv
\int_>\frac{d\varpi\,d\tilde p}{(2\pi)^2}\,
\frac{1}{\tilde p^{\,2}+\tilde\rho_r(\varpi)\varpi^2}
\simeq
\frac{d\ell}{2\pi\sqrt{\tilde\rho_{r,\ell}}},
\end{equation}
with
\begin{equation}
\tilde\rho_{r,\ell}
=
\tilde\rho_{0,r}+\kappa_r\ell,
\qquad
|\varpi|\sim \Omega_{0,r}e^{-\ell}.
\end{equation}
Therefore the couplings obey
\begin{equation}\label{eq:generic-r-rg-flows}
\frac{d\kappa_r}{d\ell}=0,
\qquad
\frac{d\tilde\rho_{0,r}}{d\ell}=\kappa_r,
\qquad
\frac{d\lambda_r}{d\ell}
=
-\frac{R}{\pi\sqrt{\tilde\rho_{r,\ell}}}
+O(\lambda_r^{-1}).
\end{equation}
Thus the one-loop beta function has exactly the same algebraic form as in the
\(r=1\) analysis.  The qualitative difference is purely kinematic: the shell is
centered on the ballistic line \(\Omega=c_r k\), not on \(\Omega=0\).

For any reference scale \(\ell_0\) at which \(\tilde\rho_{r,\ell_0}>0\), the
flow integrates to
\begin{equation}\label{eq:generic-r-lambda-integrated}
\lambda_r(\ell)
=
\lambda_r(\ell_0)
-\frac{2R}{\pi\kappa_r}
\left[
\sqrt{\tilde\rho_{r,\ell}}
-\sqrt{\tilde\rho_{r,\ell_0}}
\right].
\end{equation}
So once the logarithm dominates the analytic offset, the stiffness again
decreases only as \(\lambda_r(\ell)\sim\lambda_r(\ell_0)-{\rm const}\times
\sqrt{\ell}\).  The generic-\(r\) anisotropy therefore changes the bare
stiffness, the effective velocity, and the direction of the soft mode, but it
does not change the conclusion that the logarithmic kernel slows the flow
toward disorder and enhances the correlation length parametrically.

\subsection{Running of the aspect ratio}

It is useful to separate three distinct notions of anisotropy.  First, the
ballistic velocity \(c_r\) selects the singular direction
\(\delta_r=\Omega-c_rk\).  This is fixed by the endpoint kinematics of the
fermion bubble and is not a coupling of the sigma model itself.  Second, the
analytic mixed term in the ballistic frame is encoded in the dimensionless ratio
\(v_r\).  Because the one-loop shell integration renormalizes the whole bracket
in \eqref{eq:SE-Y-generic-r-kernel} by the same factor, \(v_r\) is not
renormalized at this order:
\begin{equation}
\frac{dv_r}{d\ell}=0+O(\lambda_r^{-1}).
\end{equation}
Thus the orientation of the soft ridge is fixed; in the original laboratory
variables the shell remains centered on the ballistic line \(\Omega=c_rk\).

The actual running anisotropy is the relative weight of \(\tilde p^2\) and
\(\varpi^2\) after both shears have been performed.  At shell scale \(\ell\), the
quadratic kernel is
\[
\tilde p^{\,2}+\tilde\rho_{r,\ell}\varpi^2,
\qquad
\tilde\rho_{r,\ell}=\tilde\rho_{0,r}+\kappa_r\ell.
\]
If one defines the momentum-space shell aspect ratio by
\begin{equation}\label{eq:aspect-ratio-def}
a_r(\ell)\equiv \sqrt{\tilde\rho_{r,\ell}},
\end{equation}
then typical shell momenta satisfy
\begin{equation}\label{eq:aspect-ratio-shell}
|\tilde p|_{\ell}\sim \Lambda e^{-\ell},
\qquad
|\varpi|_{\ell}\sim \frac{\Lambda e^{-\ell}}{a_r(\ell)}.
\end{equation}
Equivalently, in real space the characteristic coarse-grained scales obey
\[
\frac{L_{x'}}{L_{\tau'}}\sim a_r(\ell)=\sqrt{\tilde\rho_{r,\ell}}.
\]
Therefore the aspect ratio does flow:
\begin{equation}\label{eq:aspect-ratio-flow}
\frac{d\log a_r}{d\ell}
=
\frac{1}{2}\frac{d\log\tilde\rho_{r,\ell}}{d\ell}
=
\frac{\kappa_r}{2\tilde\rho_{r,\ell}},
\end{equation}
so at large scales
\begin{equation}
a_r(\ell)\sim \sqrt{\kappa_r\ell}
\qquad
(\ell\gg |\tilde\rho_{0,r}|/\kappa_r).
\end{equation}
The anisotropy therefore grows only logarithmically.

This logarithmic running is not the same as flowing to a new fixed dynamical
exponent.  At any given shell one can isotropize the quadratic form by the
scale-dependent rescaling
\[
\hat\varpi = a_r(\ell)\,\varpi,
\]
but there is no single scale-independent transformation that removes
\(\tilde\rho_r(\varpi)\) at all scales.  The soft dispersion is
\begin{equation}\label{eq:generic-r-soft-dispersion}
|\varpi|_\ell
\sim
\frac{|\tilde p|_\ell}{\sqrt{\tilde\rho_{r,\ell}}}
\sim
\frac{|\tilde p|_\ell}{\sqrt{\kappa_r\ell}}
\sim
\frac{|\tilde p|_\ell}
{\sqrt{\kappa_r\log(\Lambda/|\tilde p|_\ell)}}.
\end{equation}
Thus the leading scaling remains \(z=1\), but with a multiplicative logarithmic
correction rather than a pure power law.  If one insists on defining a local
effective exponent from \(|\varpi|_\ell\propto |\tilde p|_\ell^{\,z_{\rm eff}}\),
then
\begin{equation}
z_{\rm eff}(\ell)
=
1+\frac{\kappa_r}{2\tilde\rho_{r,\ell}}
\xrightarrow[\ell\to\infty]{}1.
\end{equation}
So the correct statement is that the aspect ratio runs, but there is no
separate one-loop flow of the ballistic velocity or of a universal anisotropic
critical exponent.

\appendix

\section{Fermion-Bubble Input to the NLSM}
\label{app:bubble-appendix}

\subsection{Retarded-advanced bubble}
\label{app:ra-bubble}

The retarded Green's function on the semi-infinite geometry is
\begin{equation}
G_R^{-1}(\omega,q)
=
(\omega+i\Gamma/2)\,\mathbf{1}
-\sin q\,\sigma_x
-(1-\cos q)\,\sigma_z
-\Sigma_R(\omega,q)\,P_-,
\end{equation}
with
\begin{equation}
\Sigma_R(\omega,q)
=
\frac{
\omega^2+2\cos q-1
-\sqrt{\omega^2-1+i0}\,\sqrt{\omega^2-(5-4\cos q)+i0}
}{
2(\omega-\sin q+i0)
}.
\end{equation}
The retarded-advanced bubble entering the quadratic NLSM is
\begin{equation}
\Pi_{RA}(\Omega,k)
=
\int\frac{d\omega\,dq}{(2\pi)^2}\,
\tr_{\rm orb}\!\left[
G_R(\omega+\Omega,q+k)\,G_A(\omega,q)
\right].
\end{equation}
\begin{equation}\label{eq:Delta-appendix-main}
\Delta(\Omega,k)
=
\Gamma\,\Re\!\left[
\frac{\Pi_{RA}(\Omega,k)+\Pi_{RA}(-\Omega,-k)}{2}-\Pi_{RA}(0,0)
\right].
\end{equation}
For the line-cut benchmarks below we use direct midpoint quadrature on the
trusted small windows
\[
5\times10^{-4}\le k\le 2.5\times10^{-3},
\qquad
5\times10^{-4}\le \Omega\le 1.5\times10^{-3},
\]
with datasets
\texttt{data/smallk\_direct\_lines\_Nw16384\_Nq4096\_wmax320} and
\texttt{data/smallomega\_direct\_lines\_Nw32768\_Nq2048\_wmax640}.  The largest
damping used in these fits is \(\Gamma=80\), so the numerical frequency
bandwidth always exceeds \(\Gamma\).

It is also useful to resolve the bubble in the basis where \(\sigma_x\) is
diagonal:
\begin{equation}
G'(\omega,q)=
\begin{pmatrix}
G_{++}(\omega,q) & G_{+-}(\omega,q)\\
G_{-+}(\omega,q) & G_{--}(\omega,q)
\end{pmatrix}.
\end{equation}
Because the Frobenius overlap is basis-invariant, the RA bubble splits into
component-wise contributions of \(G_{++}\), \(G_{+-}\), \(G_{-+}\), and
\(G_{--}\).  The exact numerical decomposition shows that the low-energy
coefficients are dominated by the \(--\) channel.

\subsubsection{Zero-momentum bubble}

At zero external arguments,
\begin{equation}\label{eq:pi00-app}
\Pi_{RA}(0,0)
=
\int\frac{d\omega}{2\pi}
\int_{-\pi}^{\pi}\frac{dq}{2\pi}
\sum_{a,b=1}^2
|G_{R,ab}(\omega,q)|^2.
\end{equation}
This was evaluated directly with adaptive bandwidth
\[
\omega_{\max}=\max(320,8\Gamma).
\]
The raw result has the universal high-frequency tail
\[
\Pi_{RA}(0,0)
=
\Pi_{\mathrm{finite}\;\omega_{\max}}(0,0)
+\frac{2}{\pi\omega_{\max}}
+o(\omega_{\max}^{-1}),
\]
so after tail correction the data collapse rapidly to
\begin{equation}\label{eq:pi00-limit-app}
\Gamma\Pi_{RA}(0,0)\to 2.
\end{equation}
Numerically, the corrected totals are
\[
1.98597,\quad 1.99545,\quad 1.99878,\quad 1.99975,\quad 2.00001
\]
for \(\Gamma=10,20,40,80,160\).

\begin{figure}[p]
    \centering
    \includegraphics[width=\textwidth]{output/pi00_direct_analysis.pdf}
    \caption{Direct numerical evaluation of \(\Gamma\Pi_{RA}(0,0)\).  Top: raw and tail-corrected totals.  Middle: decomposition in the \(\sigma_x\)-diagonal basis.  Bottom: bandwidth dependence and the corresponding tail correction.}
    \label{fig:pi00-app}
\end{figure}

The same limit follows analytically from the large-\(\Gamma\) expansion
\[
G_R(\omega,q)
=
\frac{2}{\Gamma(x+i)}\,\mathbf{1} + O(\Gamma^{-2}),
\qquad
\omega=\frac{\Gamma x}{2},
\]
which gives
\begin{equation}
\Gamma\Pi_{RA}(0,0)
=
\int\frac{dq}{2\pi}\int\frac{dx}{\pi}\frac{2}{x^2+1}
+O(\Gamma^{-1})
=
2+O(\Gamma^{-1}).
\end{equation}
It is useful to spell out this limit more explicitly.  Write
\[
G_R^{-1}(\omega,q)
=
\bigl(\omega+i\Gamma/2\bigr)\mathbf{1}-\mathcal M(\omega,q),
\]
where
\[
\mathcal M(\omega,q)=
\begin{pmatrix}
\sin q & 1-\cos q\\
1-\cos q & -\sin q+\Sigma_R(\omega,q)
\end{pmatrix}.
\]
For fixed \(q\) and scaled frequency \(\omega=\Gamma x/2\), the exact bulk
self-energy satisfies \(\Sigma_R(\Gamma x/2,q)=O(\Gamma^{-1})\) at fixed \(x\),
so \(\mathcal M(\Gamma x/2,q)=O(1)\) uniformly in \(q\).  Therefore
\[
G_R(\Gamma x/2,q)
=
\frac{2}{\Gamma}
\left[
(x+i)\mathbf{1}-\frac{2}{\Gamma}\mathcal M(\Gamma x/2,q)
\right]^{-1}
=
\frac{2}{\Gamma(x+i)}\,\mathbf{1}+O(\Gamma^{-2}).
\]
Substituting into \eqref{eq:pi00-app} yields
\begin{align}
\Gamma\Pi_{RA}(0,0)
&=
\Gamma\int\frac{dq}{2\pi}\int\frac{d\omega}{2\pi}\,
\mathrm{Tr}\!\left[G_R(\omega,q)G_R(\omega,q)^\dagger\right]
\notag\\
&=
\Gamma\int\frac{dq}{2\pi}\int\frac{\Gamma\,dx}{4\pi}\,
\mathrm{Tr}\!\left[
\frac{4}{\Gamma^2(x^2+1)}\,\mathbf{1}+O(\Gamma^{-3})
\right]
\notag\\
&=
\int\frac{dq}{2\pi}\int\frac{dx}{\pi}\,\frac{2}{x^2+1}
+O(\Gamma^{-1})
=
2+O(\Gamma^{-1}).
\end{align}

The \(O(\Gamma^{-1})\) remainder above is only a bound from the broad
frequency-scaling expansion.  The defect in the exact sum rule is smaller.  The
retarded and advanced Green's functions obey
\[
G_R-G_A
=
G_R\Bigl[(\Sigma_R-\Sigma_A)P_- - i\Gamma\,\mathbf{1}\Bigr]G_A,
\]
so
\[
i(G_R-G_A)
=
\Gamma G_RG_A
+i\,G_R(\Sigma_R-\Sigma_A)P_-G_A.
\]
After integrating over frequency and momentum, the spectral sum rule on the
left gives \(2\).  Therefore
\begin{equation}
\label{eq:pi00-mass-identity-app}
2-\Gamma\Pi_{RA}(0,0)
=
\int\frac{d\omega\,dq}{(2\pi)^2}\,
\bigl[-2\,\mathrm{Im}\,\Sigma_R(\omega,q)\bigr]\,
\tr_{\rm orb}\!\left[
G_R(\omega,q)P_-G_A(\omega,q)
\right].
\end{equation}
The factor \(-2\,\mathrm{Im}\,\Sigma_R\) is supported only on the finite bulk
continuum.  On that support, \(G_R=O(\Gamma^{-1})\), and hence
\begin{equation}
\label{eq:pi00-mass-scale-app}
2-\Gamma\Pi_{RA}(0,0)=O(\Gamma^{-2}).
\end{equation}

There is also a stronger exact statement for the pole-only approximation.  Let
\[
Z_+(q)=
\begin{cases}
\cos q\,(2-\cos q), & \cos q>0,\\
0, & \cos q<0,
\end{cases}
\qquad
\Sigma_{\mathrm{pole}}(z,q)=\frac{Z_+(q)}{z-\sin q},
\]
and denote by \(G_R^{\mathrm{pole}}(z,q)\) the Green's function obtained by
replacing \(\Sigma_R\) with \(\Sigma_{\mathrm{pole}}\).  This pole-only
Green's function is the physical \(2\times2\) block of the Hermitian
\(3\times3\) resolvent
\[
\mathcal G^{\mathrm{pole}}(z,q)
=
\left[
z\mathbf{1}_3
-
\begin{pmatrix}
\sin q & 1-\cos q & 0\\
1-\cos q & -\sin q & \sqrt{Z_+(q)}\\
0 & \sqrt{Z_+(q)} & \sin q
\end{pmatrix}
+
\frac{i\Gamma}{2}
\begin{pmatrix}
1&0&0\\
0&1&0\\
0&0&0
\end{pmatrix}
\right]^{-1},
\]
namely \(G_R^{\mathrm{pole}}=P\,\mathcal G^{\mathrm{pole}}P\) with
\[
P=
\begin{pmatrix}
1&0&0\\
0&1&0\\
0&0&0
\end{pmatrix}.
\]
Because the \(3\times3\) Hamiltonian is Hermitian, its resolvent obeys
\[
i\bigl(\mathcal G^{\mathrm{pole}}-\mathcal G^{\mathrm{pole}\,\dagger}\bigr)
=
\Gamma\,\mathcal G^{\mathrm{pole}}P\mathcal G^{\mathrm{pole}\,\dagger}.
\]
Projecting onto the physical block gives
\[
i\bigl(G_R^{\mathrm{pole}}-G_R^{\mathrm{pole}\,\dagger}\bigr)
=
\Gamma\,G_R^{\mathrm{pole}}G_R^{\mathrm{pole}\,\dagger}.
\]
Therefore
\[
\Gamma\,\mathrm{Tr}\!\left[
G_R^{\mathrm{pole}}(\omega,q)G_R^{\mathrm{pole}}(\omega,q)^\dagger
\right]
=
-2\,\mathrm{Im}\,\mathrm{Tr}\,G_R^{\mathrm{pole}}(\omega,q).
\]
Since \(G_R^{\mathrm{pole}}(z,q)=\mathbf{1}_2/z+O(z^{-2})\) as \(|z|\to\infty\),
the spectral sum rule gives
\[
\int_{-\infty}^{\infty}\frac{d\omega}{\pi}\,
\bigl[-\mathrm{Im}\,\mathrm{Tr}\,G_R^{\mathrm{pole}}(\omega,q)\bigr]=2
\]
for every \(q\).  Hence the pole-only zero-momentum bubble is exact:
\begin{equation}
\Gamma\Pi_{RA}^{\mathrm{pole}}(0,0)=2.
\end{equation}

\subsubsection{Purely imaginary linear terms}
\label{app:ra-bubble-odd-linear}

The appendix results used so far concern the symmetrized real part
\eqref{eq:Delta-appendix-main}.  For the \((\theta,\phi)\) sector one also needs
the odd analytic part of the full complex bubble.  Expanding the retarded-
advanced bubble at small external arguments gives
\begin{equation}\label{eq:Pi-linear-expansion-app}
\Pi_{RA}(\Omega,k)-\Pi_{RA}(0,0)
=
\Omega\,\mathcal I_\Omega
+k\,\mathcal I_k
+O(\Omega^2,\Omega k,k^2),
\end{equation}
with
\begin{equation}\label{eq:Iomega-Ik-def-app}
\mathcal I_\Omega
\equiv
\int\frac{d\omega\,dq}{(2\pi)^2}\,
\tr_{\rm orb}\!\left[(\partial_\omega G_R)(\omega,q)\,G_A(\omega,q)\right],
\qquad
\mathcal I_k
\equiv
\int\frac{d\omega\,dq}{(2\pi)^2}\,
\tr_{\rm orb}\!\left[(\partial_q G_R)(\omega,q)\,G_A(\omega,q)\right].
\end{equation}
Because \(G_A(\omega,q)=G_R(\omega,q)^\dagger\) on the real axis,
\begin{align}
\Re\,\mathcal I_\Omega
&=
\frac{1}{2}
\int\frac{d\omega\,dq}{(2\pi)^2}\,
\partial_\omega\,
\tr_{\rm orb}\!\left[G_R(\omega,q)G_A(\omega,q)\right]
=0,
\label{eq:Iomega-real-zero-app}
\\
\Re\,\mathcal I_k
&=
\frac{1}{2}
\int\frac{d\omega\,dq}{(2\pi)^2}\,
\partial_q\,
\tr_{\rm orb}\!\left[G_R(\omega,q)G_A(\omega,q)\right]
=0,
\label{eq:Ik-real-zero-app}
\end{align}
where the \(\omega\) total derivative vanishes because \(G_R\sim \omega^{-1}\)
at large frequency, while the \(q\) total derivative vanishes by periodicity.
Therefore the linear terms are necessarily pure imaginary:
\begin{equation}\label{eq:Pi-linear-imag-app}
\Pi_{RA}(\Omega,k)-\Pi_{RA}(0,0)
=
i\,\beta_\Omega(\Gamma)\,\Omega
-i\,\beta_k(\Gamma)\,k
+O(\Omega^2,\Omega k,k^2),
\end{equation}
with real coefficients \(\beta_\Omega(\Gamma)\) and \(\beta_k(\Gamma)\).

The large-\(\Gamma\) scaling of the \(\Omega\) term can be obtained directly
from the leading broadening-dominated Green's function.  Writing
\(\omega=\Gamma x/2\), one has
\[
G_R(\omega,q)=\frac{2}{\Gamma(x+i)}\,\mathbf{1}+O(\Gamma^{-2}),
\qquad
G_A(\omega,q)=\frac{2}{\Gamma(x-i)}\,\mathbf{1}+O(\Gamma^{-2}).
\]
Keeping only this universal leading piece gives
\begin{equation}\label{eq:PiOmega-leading-app}
\Pi_{RA}(\Omega,0)
=
\frac{2}{\Gamma}
\int\frac{dx}{\pi}\,
\frac{1}{(x+2\Omega/\Gamma+i)(x-i)}
+O(\Gamma^{-2})
=
\frac{2}{\Gamma}\,
\frac{1}{1-i\Omega/\Gamma}
+O(\Gamma^{-2}),
\end{equation}
hence
\begin{equation}\label{eq:Omega-linear-leading-app}
\Gamma\bigl[\Pi_{RA}(\Omega,0)-\Pi_{RA}(0,0)\bigr]
=
\frac{2i\,\Omega/\Gamma}{1-i\Omega/\Gamma}
+O(\Gamma^{-2})
=
\frac{2i\,\Omega}{\Gamma}
+O\!\left(\frac{\Omega^2}{\Gamma^2},\Gamma^{-2}\right).
\end{equation}
So the odd frequency term is analytic and scales as \(1/\Gamma\) in
\(\Gamma\Pi_{RA}\).

For the momentum term, the same leading broadening-only approximation gives no
\(k\) dependence at all, because the leading Green's function is proportional to
the identity and independent of \(q\).  The linear \(k\) term therefore first
appears through the subleading \(q\)-dependent corrections to \(G_R\), so
\begin{equation}\label{eq:k-linear-scaling-app}
\Gamma\,\Im\!\left[\Pi_{RA}(0,k)-\Pi_{RA}(0,0)\right]
=
O\!\left(\frac{k}{\Gamma}\right).
\end{equation}
Thus both odd linear terms decay with \(\Gamma\), although the \(k\) term is
absent in the leading broadening approximation.

A direct recomputation of the full complex line cuts using
\texttt{src/collect\_small\_argument\_direct\_lines.jl} on a moderate grid
\[
N_\omega=4096,\qquad N_q=1024,\qquad \omega_{\max}=320,\qquad \eta=10^{-4},
\]
confirms this scaling.  At \(\Omega=\pm 5\times10^{-4},\pm 10^{-3}\) and
\(k=\pm 5\times10^{-4},\pm 10^{-3}\), one finds
\begin{equation}\label{eq:odd-linear-numerics-app}
\frac{\Gamma\,\Im[\Pi_{RA}(\Omega,0)-\Pi_{RA}(0,0)]}{\Omega}
\approx
\frac{2.379}{\Gamma},\quad
\frac{2.377}{\Gamma},\quad
\frac{2.374}{\Gamma},
\end{equation}
\begin{equation}\label{eq:k-linear-numerics-app}
\frac{\Gamma\,\Im[\Pi_{RA}(0,k)-\Pi_{RA}(0,0)]}{k}
\approx
-\frac{0.313}{\Gamma},\quad
-\frac{0.304}{\Gamma},\quad
-\frac{0.299}{\Gamma},
\end{equation}
for \(\Gamma=20,40,80\), respectively.  The numerical prefactors are therefore
small and decrease approximately as \(1/\Gamma\), in agreement with the
large-\(\Gamma\) analytic estimate above.

\subsubsection{Small-\texorpdfstring{$k$}{k} behavior}

On the trusted small-\(k\) window, the line cut is already very close to a pure
quadratic:
\[
\Re[\Gamma(\Pi_{RA}(0,k)-\Pi_{RA}(0,0))]
\approx
-A_k(\Gamma)\,k^2.
\]
The fitted coefficients are
\[
A_k(20)\approx 0.07322,\qquad
A_k(40)\approx 0.06882,\qquad
A_k(80)\approx 0.06730,
\]
with sub-\(10^{-3}\) relative RMS errors.  Figure~\ref{fig:Ak-app} shows the
channel decomposition.

\begin{figure}[p]
    \centering
    \includegraphics[width=\textwidth]{output/edge_Ak_model_comparison.pdf}
    \caption{Decomposition of the small-\(k\) line in the \(\sigma_x\)-diagonal basis.  The pole-only \(++\) edge estimate fails badly, while the exact \(--\) channel nearly saturates the full direct result.}
    \label{fig:Ak-app}
\end{figure}

Numerically, the corresponding quadratic coefficients are
\begin{center}
\begin{tabular}{crrrrr}
\toprule
\(\Gamma\) & full & pole-only edge & exact \(++\) & off-diagonal & exact \(--\) \\
\midrule
20 & 0.07322 & 1.03780 & 0.00126 & 0.00249 & 0.06947 \\
40 & 0.06882 & 1.03729 & 0.00031 & 0.00064 & 0.06786 \\
80 & 0.06730 & 1.03717 & 0.00008 & 0.00016 & 0.06706 \\
\bottomrule
\end{tabular}
\end{center}

To derive the coefficient analytically, define for any single component
\begin{equation}
\Pi_g(0,k)=
\int\frac{d\omega\,dq}{(2\pi)^2}\,
g(\omega,q)\,\overline{g(\omega,q-k)}.
\end{equation}
Then
\begin{equation}\label{eq:Ak-gradient-app}
\frac{\Pi_g(0,k)+\Pi_g(0,-k)}{2}-\Pi_g(0,0)
=
-\frac{k^2}{2}\int\frac{d\omega\,dq}{(2\pi)^2}\,
|\partial_q g(\omega,q)|^2
+O(k^4),
\end{equation}
so
\begin{equation}
A_k^{(g)}(\Gamma)
=
\frac{\Gamma}{2}\int\frac{d\omega\,dq}{(2\pi)^2}
|\partial_q g(\omega,q)|^2.
\end{equation}
For the dominant \(g=G_{--}\) channel, define
\[
Z(q)=\cos q\,(2-\cos q).
\]
Inside the edge interval \(\cos q>0\), the bulk self-energy has the pole
\[
\Sigma_R(\omega,q)=\frac{Z(q)}{\omega-\sin q}+\Sigma_{\rm reg}(\omega,q).
\]
Here ``edge interval'' means \(q\in(-\pi/2,\pi/2)\).  For \(\cos q<0\), the
numerator of the exact self-energy vanishes at \(\omega=\sin q\), so the
would-be \(1/x\) pole is absent and the leading resonance-strip contribution
does not arise.  From the matrix in the \(\sigma_x\)-diagonal basis,
\[
G_R^{-1}(\omega,q)=
\begin{pmatrix}
A & -m\\
-m & B
\end{pmatrix},
\qquad
A=\omega-\sin q+i\Gamma/2,
\qquad
m=1-\cos q,
\]
with
\[
B=\omega+\sin q-\Sigma_R(\omega,q)+i\Gamma/2,
\]
we have
\[
G_{--}(\omega,q)=\frac{1}{B-m^2/A}.
\]
At large \(\Gamma\), \(m^2/A=O(\Gamma^{-1})\), so
\[
G_{--}(\omega,q)
=
\frac{1}{\omega+\sin q-\Sigma_R(\omega,q)+i\Gamma/2}
+O(\Gamma^{-2}),
\]
which is the starting point of the asymptotic analysis.

The full derivative at fixed \(\omega\) is
\[
\partial_q G_{--}
=
\frac{\partial_q\Sigma_R(\omega,q)-\cos q}
{\bigl(\omega+\sin q-\Sigma_R(\omega,q)+i\Gamma/2\bigr)^2}
+O(\Gamma^{-1}),
\]
and for the pole part of the self-energy,
\[
\partial_q\Sigma_R
=
\frac{Z'(q)}{\omega-\sin q}
+
\frac{Z(q)\cos q}{(\omega-\sin q)^2}
+
\partial_q\Sigma_{\rm reg}.
\]
On the resonance strip
\[
\omega-\sin q = O\!\left(\frac{Z(q)}{\Gamma}\right),
\qquad
\omega+\sin q-\Sigma_R(\omega,q)+i\Gamma/2 = O(\Gamma).
\]
The four numerator terms then scale as
\[
\cos q=O(1),\qquad
\partial_q\Sigma_{\rm reg}=O(1),\qquad
\frac{Z'(q)}{\omega-\sin q}=O(\Gamma),\qquad
\frac{Z(q)\cos q}{(\omega-\sin q)^2}=O(\Gamma^2).
\]
After division by the \(O(\Gamma^2)\) denominator in \(\partial_q G_{--}\),
only the last term survives at leading order.  This is why the explicit
derivative of \(+\sin q\) and the regular part of the self-energy are
subleading in the large-\(\Gamma\) asymptotics.

Therefore the leading contribution to the gradient formula can be written
directly in the original variables as
\begin{equation}
\label{eq:Ak-omegaq-app-detailed}
A_k^{(--)}(\Gamma)
\sim
\frac{\Gamma}{2}
\int_{-\pi/2}^{\pi/2}\frac{dq}{2\pi}
\int\frac{d\omega}{2\pi}\;
\frac{Z(q)^2\cos^2 q}
{(\omega-\sin q)^4\,
\left|\omega+\sin q-\frac{Z(q)}{\omega-\sin q}+i\Gamma/2\right|^4}.
\end{equation}
This integral is concentrated on the narrow strip
\[
\omega-\sin q = O\!\left(\frac{Z(q)}{\Gamma}\right),
\]
so it is convenient only at this stage to set
\[
\omega-\sin q=\frac{Z(q)}{\Gamma}u.
\]
Then
\[
\omega+\sin q-\frac{Z(q)}{\omega-\sin q}+i\Gamma/2
=
\Gamma\left(\frac{i}{2}-\frac{1}{u}\right)+O(1),
\]
and
\[
\partial_q G_{--}
=
\frac{4\cos q}{Z(q)(iu-2)^2}+O(\Gamma^{-1}).
\]
Substituting back into the gradient formula yields
\begin{equation}\label{eq:Ak-final-app}
A_k^{(--)}(\infty)
=
\frac{1}{8\pi}\int_{-\pi/2}^{\pi/2}\frac{\cos q}{2-\cos q}\,dq
=
\frac{1}{3\sqrt{3}}-\frac{1}{8}
\approx 0.06745009.
\end{equation}
This agrees very well with the extracted exact-\(--\) coefficients above, so the
full small-\(k\) coefficient inherits the same limit.

\subsubsection{Small-\texorpdfstring{$\Omega$}{Omega} behavior}

On the trusted small-\(\Omega\) window,
\[
\Re[\Gamma(\Pi_{RA}(\Omega,0)-\Pi_{RA}(0,0))]
\approx
-\widehat A_\Omega(\Gamma)\,\Omega^2.
\]
The fitted full coefficients are
\[
\widehat A_\Omega(20)\approx 0.38406,\qquad
\widehat A_\Omega(40)\approx 0.35505,\qquad
\widehat A_\Omega(80)\approx 0.33140.
\]
Again the exact \(--\) channel dominates, as shown in
Fig.~\ref{fig:Aomega-app}.

\begin{figure}[p]
    \centering
    \includegraphics[width=\textwidth]{output/edge_Aomega_model_comparison.pdf}
    \caption{Decomposition of the small-\(\Omega\) line in the \(\sigma_x\)-diagonal basis.  The exact \(--\) contribution already saturates the full direct line for \(\Gamma=40\) and \(80\).}
    \label{fig:Aomega-app}
\end{figure}

The fitted coefficients are
\begin{center}
\begin{tabular}{crrrrr}
\toprule
\(\Gamma\) & full & exact \(++\) & exact \(+-\) & exact \(-+\) & exact \(--\) \\
\midrule
20 & 0.38406 & -0.00173 & 0.00419 & 0.00419 & 0.37742 \\
40 & 0.35505 & 0.00032 & 0.00066 & 0.00066 & 0.35341 \\
80 & 0.33140 & 0.00014 & 0.00013 & 0.00013 & 0.33101 \\
\bottomrule
\end{tabular}
\end{center}

For a single component,
\begin{equation}
\frac{\Pi_g(\Omega,0)+\Pi_g(-\Omega,0)}{2}-\Pi_g(0,0)
=
-\frac{\Omega^2}{2}\int\frac{d\omega\,dq}{(2\pi)^2}\,
|\partial_\omega g(\omega,q)|^2
+O(\Omega^4),
\end{equation}
so
\begin{equation}
A_\Omega^{(g)}(\Gamma)
=
\frac{\Gamma}{2}\int\frac{d\omega\,dq}{(2\pi)^2}
|\partial_\omega g(\omega,q)|^2.
\end{equation}
For \(g=G_{--}\), the same large-\(\Gamma\) form
\[
G_{--}(\omega,q)
=
\frac{1}{\omega+\sin q-\Sigma_R(\omega,q)+i\Gamma/2}
+O(\Gamma^{-2})
\]
still applies.  On the resonance strip
\[
\omega-\sin q=\frac{Z(q)}{\Gamma}u,
\]
the pole part of the self-energy gives
\[
\partial_\omega\Sigma_R \sim -\frac{Z(q)}{(\omega-\sin q)^2},
\qquad
\partial_\omega G_{--}
=
\frac{4}{Z(q)(iu-2)^2}+O(\Gamma^{-1}).
\]
Substituting into the gradient formula gives the formal density
\[
A_\Omega^{(--)}(\Gamma)\sim \frac{1}{8\pi}\int_{-\pi/2}^{\pi/2}\frac{dq}{Z(q)},
\]
which diverges logarithmically at the touching points because \(Z(q)\to 0\) as
\(q\to\pm\pi/2\).  Keeping \(\Omega\neq 0\) as the infrared regulator yields
\begin{equation}\label{eq:Aomega-final-app}
\Re[\Gamma(\Pi_{RA,--}(\Omega,0)-\Pi_{RA,--}(0,0))]_{\rm sing}
=
-\frac{\Omega^2}{8\pi}\log\!\frac{C_\Omega}{\Gamma|\Omega|}
+O(\Omega^2).
\end{equation}
This does not directly predict the finite-window fitted coefficient
\(\widehat A_\Omega(\Gamma)\), but it does imply the drift
\[
\widehat A_{\Omega,--}(\Gamma)
=
B_\Omega-\frac{1}{8\pi}\log\Gamma+o(1).
\]
Numerically, the shifted combination
\[
\widehat A_{\Omega,--}(\Gamma)+\frac{1}{8\pi}\log\Gamma
\]
is nearly constant:
\[
0.49661,\qquad 0.50018,\qquad 0.50536
\qquad
(\Gamma=20,40,80),
\]
which is the appropriate benchmark of the logarithmic prediction.

\subsubsection{Mixed \texorpdfstring{$\Omega k$}{Omega k} term}

To isolate the bilinear term, define the odd-odd projection
\begin{equation}\label{eq:Aomegak-oddodd-app}
\Delta_{\mathrm{odd\text{-}odd}}(\Omega,k)
=
\frac{
\Gamma\,\Re\Pi_{RA}(\Omega,k)
-\Gamma\,\Re\Pi_{RA}(\Omega,-k)
-\Gamma\,\Re\Pi_{RA}(-\Omega,k)
+\Gamma\,\Re\Pi_{RA}(-\Omega,-k)
}{4}.
\end{equation}
For small arguments,
\[
\Delta_{\mathrm{odd\text{-}odd}}(\Omega,k)
=
A_{\Omega k}(\Gamma)\,\Omega k+o(\Omega k).
\]
The direct decomposition shows that this coefficient is also overwhelmingly the
exact \(--\) channel.  At the small point
\[
\Omega=k=5\times10^{-4},
\qquad
N_\omega=16384,
\qquad
N_q=4096,
\qquad
\omega_{\max}=640,
\]
the odd-odd quotient gives
\begin{center}
\begin{tabular}{crrrr}
\toprule
\(\Gamma\) & full & exact \(++\) & off-diagonal & exact \(--\) \\
\midrule
20 & 0.19760 & 0.00011 & 0.00074 & 0.19675 \\
40 & 0.19283 & 0.00001 & 0.00016 & 0.19266 \\
80 & 0.19038 & 0.00000 & 0.00004 & 0.19034 \\
\bottomrule
\end{tabular}
\end{center}
At the still smaller point \(\Omega=k=2.5\times10^{-4}\), one finds
\[
A_{\Omega k}^{\rm full}(80)\approx 0.19130,
\qquad
A_{\Omega k}^{(--)}(80)\approx 0.19126.
\]

For a single component,
\begin{equation}
A_{\Omega k}^{(g)}(\Gamma)
=
-\Gamma\,\Re\!\int\frac{d\omega\,dq}{(2\pi)^2}\,
(\partial_\omega g)(\partial_q \overline g).
\end{equation}
Applying this to \(g=G_{--}\), and using
\[
\partial_\omega G_{--}
=
-\frac{4}{Z(q)(iu-2)^2}+O(\Gamma^{-1}),
\qquad
\partial_q G_{--}
=
\frac{4\cos q}{Z(q)(iu-2)^2}+O(\Gamma^{-1}),
\]
on the same resonance strip, one obtains
\begin{equation}\label{eq:Aomegak-final-app}
A_{\Omega k}^{(--)}(\infty)
=
\frac{1}{4\pi}\int_{-\pi/2}^{\pi/2}\frac{dq}{2-\cos q}
=
\frac{1}{3\sqrt3}
\approx 0.19245009.
\end{equation}
This matches the direct high-resolution values above very well.

Collecting the results,
\begin{equation}\label{eq:Delta-summary-app}
\Delta(\Omega,k)
=
-A_k\,k^2
+A_{\Omega k}\,\Omega k
-\frac{\Omega^2}{8\pi}\log\!\frac{C_\Omega}{\Gamma|\Omega|}
+\cdots,
\end{equation}
with
\begin{equation}
A_k\xrightarrow[\Gamma\to\infty]{}\frac{1}{3\sqrt3}-\frac18,
\qquad
A_{\Omega k}\xrightarrow[\Gamma\to\infty]{}\frac{1}{3\sqrt3},
\qquad
\Gamma\Pi_{RA}(0,0)\xrightarrow[\Gamma\to\infty]{}2.
\end{equation}
These are precisely the coefficients used in the main text.

\subsubsection{Reduced \texorpdfstring{$G_{--}$}{G--} bubble and direct kernel expansion}

Within the leading large-\(\Gamma\) resonance-strip approximation, the
\(--\) Green's function can be replaced by the local model
\begin{equation}
\label{eq:Gmm-local-app}
G_{--}^{\mathrm{loc}}(\omega,q)
=
\frac{\omega-\sin q}{i\Gamma(\omega-\sin q)/2-Z(q)},
\qquad
Z(q)=\cos q\,(2-\cos q),
\end{equation}
for \(q\in(-\pi/2,\pi/2)\).  This is not the full lattice Green's function; it
is the leading edge-strip model that keeps only the pole part of the bulk
self-energy.

Whenever both \(q\) and \(q-k\) remain inside the edge interval, the local
\(--\) bubble is
\[
\Pi_{RA,--}^{\mathrm{loc}}(\Omega,k)
=
\int_{-\pi/2}^{\pi/2}\frac{dq}{2\pi}
\int\frac{d\omega}{2\pi}\,
G_{--}^{\mathrm{loc}}(\omega+\Omega,q+k)\,
\overline{G_{--}^{\mathrm{loc}}(\omega,q)}.
\]
The \(\omega\) integral can be done explicitly by residues.  Closing the
contour in the upper half-plane picks the pole of the advanced factor and gives
\begin{equation}
\label{eq:Gmm-local-kernel-app}
\Pi_{RA,--}^{\mathrm{loc}}(\Omega,k)
=
\int_{-\pi/2}^{\pi/2}\frac{dq}{2\pi}\;
\mathcal I_\Gamma(q;\Omega,k),
\end{equation}
with
\begin{equation}
\label{eq:Gmm-local-kernel-density-app}
\mathcal I_\Gamma(q;\Omega,k)
=
\frac{4 Z(q)}{\Gamma^2}\,
\frac{i\bigl(\Omega+\sin q-\sin(q+k)\bigr)-2Z(q)/\Gamma}
{Z(q)+Z(q+k)-i\Gamma\bigl(\Omega+\sin q-\sin(q+k)\bigr)/2}.
\end{equation}
This is a useful one-dimensional summary of the reduced \((--)\) bubble.

To expand it directly, first take its real part.  One finds
\begin{equation}
\label{eq:Gmm-local-expansion-app}
\Gamma\,\Re \mathcal I_\Gamma(q;\Omega,k)
=
-\frac{
2 Z(q)\bigl(\Omega+\sin q-\sin(q+k)\bigr)^2
\;+\;
\dfrac{8}{\Gamma^2}\,Z(q)^2\bigl(Z(q)+Z(q+k)\bigr)}
{\bigl(Z(q)+Z(q+k)\bigr)^2+\dfrac{\Gamma^2}{4}\bigl(\Omega+\sin q-\sin(q+k)\bigr)^2}.
\end{equation}
For small arguments,
\[
Z(q+k)=Z(q)+O(k),
\qquad
\Omega+\sin q-\sin(q+k)=\Omega-k\cos q+O(k^2),
\]
so
\[
\Gamma\,\Re \mathcal I_\Gamma(q;\Omega,k)
=
-\frac{4Z(q)}{\Gamma^2}
-\frac{\bigl(\Omega-k\cos q\bigr)^2}{4Z(q)}
+ O\!\left(\frac{k}{\Gamma^2},\frac{\Omega^2+k^2}{\Gamma^2},
|\Omega|^3,|\Omega|^2|k|,|\Omega|k^2,|k|^3\right).
\]
Therefore,
\begin{align}
\frac{\Gamma\,\Re \mathcal I_\Gamma(q;0,k)+
\Gamma\,\Re \mathcal I_\Gamma(q;0,-k)}{2}
-\Gamma\,\Re \mathcal I_\Gamma(q;0,0)
&=
-\frac{\cos^2 q}{4Z(q)}\,k^2
+O\!\left(k^4,\frac{k^2}{\Gamma^2}\right),
\notag\\
\frac{
\Gamma\,\Re \mathcal I_\Gamma(q;\Omega,k)
-\Gamma\,\Re \mathcal I_\Gamma(q;\Omega,-k)
-\Gamma\,\Re \mathcal I_\Gamma(q;-\Omega,k)
+\Gamma\,\Re \mathcal I_\Gamma(q;-\Omega,-k)}{4}
&=
\frac{\cos q}{2Z(q)}\,\Omega k
+o(\Omega k),
\notag\\
\frac{\Gamma\,\Re \mathcal I_\Gamma(q;\Omega,0)+
\Gamma\,\Re \mathcal I_\Gamma(q;-\Omega,0)}{2}
-\Gamma\,\Re \mathcal I_\Gamma(q;0,0)
&=
-\frac{\Omega^2}{4Z(q)}
+O\!\left(\Omega^4,\frac{\Omega^2}{\Gamma^2}\right).
\label{eq:Gmm-local-densities-app}
\end{align}
After integrating over \(q\) with measure \(dq/(2\pi)\), these direct kernel
expansions reproduce exactly the coefficient densities used above:
\[
\frac{\cos^2 q}{8\pi Z(q)}
=
\frac{\cos q}{8\pi(2-\cos q)},
\qquad
\frac{\cos q}{4\pi Z(q)}
=
\frac{1}{4\pi(2-\cos q)},
\qquad
\frac{1}{8\pi Z(q)}.
\]
The first two integrate to \eqref{eq:Ak-final-app} and
\eqref{eq:Aomegak-final-app}.  The third is not integrable at
\(q=\pm\pi/2\), which is the direct kernel-level origin of the logarithm in
\eqref{eq:Aomega-final-app}.

\subsection{Large-\texorpdfstring{$\Gamma Q$}{Gamma Q} endpoint crossover}
\label{app:large-gamma-endpoint}

The coefficient extraction in Sec.~\ref{sec:generic-r-bubbles} is a
small-\((\Omega,k)\) expansion at fixed \(\Gamma\).  The same reduced endpoint
kernel has a different leading form when \(\Gamma\) is taken large while the
external arguments are held fixed relative to \(1/\Gamma\).  This is the origin
of the nonuniformity discussed in the main text.

For generic \(0<r<2\), write
\[
q_r=\arccos(r-1),
\qquad
c_r=r-1,
\qquad
s_r=\sqrt{r(2-r)},
\]
\[
Z_r(q)=1-(r-\cos q)^2,
\qquad
D(q)=\Omega-k\cos q.
\]
The reduced endpoint kernel is
\begin{equation}
\label{eq:largeGamma-endpoint-kernel-app}
\Delta_{\rm end}(\Omega,k)
=
-\int_{-q_r}^{q_r}\frac{dq}{2\pi}\,
\frac{Z_r(q)D(q)^2}
{4Z_r(q)^2+\Gamma^2D(q)^2/4}.
\end{equation}
The Taylor/log regime follows by expanding this expression when
\(\Gamma|D(q)|\ll Z_r(q)\) away from the endpoint, with the endpoint itself
regulated by \(\Gamma|\Omega-c_rk|\).  In contrast, when
\(\Gamma Q\gg1\) and \(D(q)\) is not parametrically small over a finite interval,
the broadening term dominates the denominator over almost all \(q\).  Then
\begin{equation}
\frac{Z_r(q)D(q)^2}
{4Z_r(q)^2+\Gamma^2D(q)^2/4}
=
\frac{4Z_r(q)}{\Gamma^2}
+\cdots,
\end{equation}
and the leading endpoint contribution saturates to
\begin{equation}
\label{eq:largeGamma-saturation-app}
\Delta_{\rm end}(\Omega,k)
=
-\frac{2}{\pi\Gamma^2}I_Z(r)+\cdots,
\qquad
I_Z(r)\equiv\int_{-q_r}^{q_r}dq\,Z_r(q).
\end{equation}
The remaining integral is elementary:
\begin{equation}
\label{eq:Iz-generic-r-app}
I_Z(r)
=
(1-2r^2)q_r+(3r+1)s_r.
\end{equation}
Thus the leading \(\Gamma Q\gg1\) endpoint term is independent of the direction
from which \((\Omega,k)\) approaches the origin, as long as the approach is
kept outside parametrically small crossover regions.  If \(D(q)\) has an
interior zero, the zero only produces a subleading correction because the
numerator also contains \(D(q)^2\).

The tangent direction \(\Omega=c_rk\) is the most delicate case because the
endpoint mismatch itself vanishes.  It is nevertheless described by the same
large-\(\Gamma Q\) saturation.  Setting \(x(q)=\cos q-c_r\) and
\(A(q)=r+1-\cos q\), so that \(Z_r(q)=x(q)A(q)\), one has
\[
D(q)=-kx(q)
\qquad
(\Omega=c_rk).
\]
Substituting this into \eqref{eq:largeGamma-endpoint-kernel-app} gives the
exact reduced endpoint expression
\begin{equation}
\label{eq:tangent-crossover-app}
\Delta_{\rm tan}(k;\Gamma)
=
-\frac{k^2}{2\pi}
\int_{-q_r}^{q_r}dq\,
\frac{x(q)A(q)}
{4A(q)^2+\Gamma^2k^2/4}.
\end{equation}
For \(\Gamma|k|\ll1\), this reduces to the regular Taylor result,
\begin{equation}
\Delta_{\rm tan}(k;\Gamma)
=
\frac{q_r-J_r}{4\pi}k^2+\cdots,
\end{equation}
which is the value obtained from
\(\Delta_{r,\rm reg}^{(--)}(c_rk,k)\).  For \(\Gamma|k|\gg1\), however,
\eqref{eq:tangent-crossover-app} gives
\begin{equation}
\label{eq:tangent-largeGamma-app}
\Delta_{\rm tan}(k;\Gamma)
=
-\frac{2}{\pi\Gamma^2}I_Z(r)
+O\!\left(\frac{1}{\Gamma^4k^2}\right).
\end{equation}
Therefore the limits do not commute:
\[
\lim_{\Gamma\to\infty}\lim_{Q\to0}
\Delta_{\rm end}(\Omega,k)
\neq
\lim_{Q\to0}\lim_{\Gamma\to\infty}
\Delta_{\rm end}(\Omega,k).
\]
The coefficient formulas in Sec.~\ref{sec:generic-r-bubbles} refer to the
first ordering, or equivalently to the regime \(\Gamma Q\ll1\).  The
saturation formula \eqref{eq:largeGamma-saturation-app} describes the opposite
regime \(\Gamma Q\gg1\).

We verified this crossover numerically using
\texttt{src/verify\_generic\_r\_regular\_terms.jl}, which directly integrates
\eqref{eq:largeGamma-endpoint-kernel-app}.  The table uses
\(\Gamma=2000\) and \(k=0.02\).  The tangent column sets
\(\Omega=c_rk\), while the off-tangent column uses \(\Omega=0.03\) and
\(k=0.02\).  The ratios compare the numerical integral to
\(\Delta_{\rm sat}=-2I_Z(r)/(\pi\Gamma^2)\).
\begin{center}
\begin{tabular}{cccc}
\toprule
\(r\) &
\(\Delta_{\rm sat}\) &
\(\Delta_{\rm tan}/\Delta_{\rm sat}\) &
\(\Delta_{\rm off}/\Delta_{\rm sat}\) \\
\midrule
0.5 & \(-5.11247\times10^{-7}\) & 0.98900 & 0.98831 \\
1.0 & \(-3.86620\times10^{-7}\) & 0.98395 & 0.98000 \\
1.5 & \(-1.74744\times10^{-7}\) & 0.97461 & 0.98717 \\
\bottomrule
\end{tabular}
\end{center}
The percent-level deviations are the expected finite-\(\Gamma Q\) corrections;
they decrease as \(\Gamma Q\) is increased.

\subsubsection{Full retarded-advanced Green's function}

The preceding computation isolates the nonuniform endpoint part of the
\(--\) channel.  For the full retarded-advanced bubble, the remaining
\(++\), \(+-\), \(-+\), and regular \(--\) pieces also contribute analytic
large-\(\Gamma\) terms.  It is useful to keep the complex, unsymmetrized
subtracted kernel
\begin{equation}
\label{eq:KRA-full-def-app}
K_{RA}(\Omega,k)
\equiv
\Gamma\left[\Pi_{RA}(\Omega,k)-\Pi_{RA}(0,0)\right].
\end{equation}
The diagonal real kernel used in \eqref{eq:Delta-appendix-main} is
\(\Re[K_{RA}(\Omega,k)+K_{RA}(-\Omega,-k)]/2\), while the odd
\(\theta\)-\(\phi\) kernel is controlled by the odd imaginary part.  Therefore
the linear terms in \(K_{RA}\) should be kept even though they drop out of
\(\Delta\).

Let
\begin{equation}
h_r(q)=\sin q\,\sigma_x+(r-\cos q)\,\sigma_z.
\end{equation}
The monitored self-energy is \(O(\Gamma^{-1})\) for
\(\omega=\Gamma x/2\), so to the order needed here
\begin{equation}
G_R\!\left(\frac{\Gamma x}{2}+\Omega,q+k\right)
=
\frac{2}{\Gamma}\frac{\mathbf{1}}{x+i+2\Omega/\Gamma}
+
\frac{4}{\Gamma^2}\frac{h_r(q+k)}{(x+i)^2}
+O(\Gamma^{-3}),
\end{equation}
and similarly for \(G_A\).  The scalar part gives exactly
\begin{equation}
\label{eq:full-scalar-largeGamma-app}
K_0(\Omega)
=
\frac{2}{1-i\Omega/\Gamma}-2
=
\frac{2i\Omega}{\Gamma}
-\frac{2\Omega^2}{\Gamma^2}
+O\!\left(\frac{\Omega^3}{\Gamma^3}\right).
\end{equation}
The \(h_r h_r\) term is the first part with momentum dependence.  At
\(\Omega=0\),
\begin{align}
K_h(0,k)
&=
\frac{2}{\Gamma^2}
\int_{-\pi}^{\pi}\frac{dq}{2\pi}\,
\tr\!\left[
h_r(q+k)h_r(q)-h_r(q)^2
\right]
+O(\Gamma^{-3})
\notag\\
&=
\frac{4(\cos k-1)}{\Gamma^2}
+O(\Gamma^{-3})
=
-\frac{2k^2}{\Gamma^2}
+O\!\left(\frac{k^4}{\Gamma^2},\Gamma^{-3}\right).
\label{eq:full-hh-largeGamma-app}
\end{align}
The trace average has no term linear in \(k\).  Since the scalar denominator is
independent of \(q\), this also implies that no \(\Omega k\) term is produced
at the displayed order.  Possible odd-in-\(k\) pieces of the full numerical
kernel are therefore subleading endpoint or self-energy corrections, not
second-order coefficients of the broadening-dominated full Green's function.

Combining the smooth full-Green contribution with the endpoint saturation
\eqref{eq:largeGamma-saturation-app}, the full complex RA kernel in the regime
\[
Q\ll1,
\qquad
\Gamma Q\gg1,
\]
is
\begin{equation}
\label{eq:full-largeGamma-KRA-app}
K_{RA}(\Omega,k)
=
\frac{2i\Omega}{\Gamma}
-
\frac{
2I_Z(r)/\pi+2\Omega^2+2k^2
}{\Gamma^2}
+R_{\rm ep}(\Omega,k;\Gamma)
+\cdots .
\end{equation}
Here \(R_{\rm ep}\) denotes the direction-dependent boundary-layer correction
left after subtracting the endpoint saturation.  If \(D(q)=\Omega-k\cos q\) is
bounded away from zero on the endpoint interval, then
\(R_{\rm ep}=O(\Gamma^{-4}Q^{-2})\).  A simple interior zero of \(D\) gives a
larger but still subleading \(O(\Gamma^{-3}Q^{-1})\) correction, and the
tangent endpoint case gives the \(O(\Gamma^{-4}k^{-2})\) correction in
\eqref{eq:tangent-largeGamma-app}.  These terms can be larger than the smooth
\(Q^2/\Gamma^2\) correction for some paths to the origin, but they are smaller
than the leading \(\Gamma^{-2}\) saturation whenever \(\Gamma Q\gg1\).

Equation \eqref{eq:full-largeGamma-KRA-app} contains the requested polynomial
terms through second order in the external arguments: a linear imaginary
\(\Omega\) term, no leading linear \(k\) term, quadratic
\(\Omega^2\) and \(k^2\) terms, and no \(\Omega k\) term at this order.

We verified \eqref{eq:full-largeGamma-KRA-app} by direct quadrature of the full
orbital trace in \eqref{eq:Pi-RA-def}, using
\path{src/verify_full_large_gammaQ.jl}.  The script uses a composite frequency
grid, dense near the edge band and coarse in the high-frequency tails.  The
table compares the numerical value of \(K_{RA}\) to the displayed asymptote.
\begin{center}
\scriptsize
\begin{tabular}{cccccccc}
\toprule
\(r\) & \(\Gamma\) & \(\Omega\) & \(k\) &
\(\Re K_{\rm num}\) & \(\Re K_{\rm asy}\) &
\(\Im K_{\rm num}\) & \(\Im K_{\rm asy}\) \\
\midrule
0.5 & 200 & 0.100 & 0     & \(-4.93431\times10^{-5}\) & \(-5.16247\times10^{-5}\) & \(1.00794\times10^{-3}\) & \(1.00000\times10^{-3}\) \\
1.0 & 200 & 0.100 & 0     & \(-3.73623\times10^{-5}\) & \(-3.91620\times10^{-5}\) & \(1.00622\times10^{-3}\) & \(1.00000\times10^{-3}\) \\
1.5 & 200 & 0.100 & 0     & \(-1.73166\times10^{-5}\) & \(-1.79744\times10^{-5}\) & \(1.00186\times10^{-3}\) & \(1.00000\times10^{-3}\) \\
1.0 & 200 & 0     & 0.100 & \(-3.63248\times10^{-5}\) & \(-3.91620\times10^{-5}\) & \(-9.23737\times10^{-6}\) & 0 \\
1.0 & 400 & 0     & 0.100 & \(-9.39823\times10^{-6}\) & \(-9.79049\times10^{-6}\) & \(-1.21641\times10^{-6}\) & 0 \\
1.5 & 200 & 0.100 & 0.050 & \(-1.67254\times10^{-5}\) & \(-1.80994\times10^{-5}\) & \(1.00337\times10^{-3}\) & \(1.00000\times10^{-3}\) \\
\bottomrule
\end{tabular}
\end{center}
The pure-\(\Omega\) rows confirm the leading imaginary term
\(2\Omega/\Gamma\) and the real \(\Omega^2/\Gamma^2\) correction.  The pure
\(k\) rows test the real \(k^2/\Gamma^2\) correction together with the endpoint
saturation.  The small residual imaginary value in those rows is odd in \(k\)
and falls by nearly a factor of eight when \(\Gamma\) is doubled from \(200\)
to \(400\), consistent with a subleading remainder rather than a leading linear
\(k\) term.  The mixed row checks that no \(\Omega k\) coefficient is needed at
this order.

\subsection{Full Keldysh-Keldysh bubble}
\label{app:kk-bubble}

For the full Green's function used in the numerics, the Keldysh component is not
obtained from a pole-only ansatz.  Instead, the upper-triangular Dyson
construction gives
\begin{equation}\label{eq:GK-dyson-appendix}
G_K(\omega,q)
=
\sgn(\omega)\Bigl[
G_R(\omega,q)-G_A(\omega,q)+i\Gamma\,G_R(\omega,q)G_A(\omega,q)
\Bigr].
\end{equation}
Using
\begin{equation}
G_R-G_A
=
G_R\Bigl[(\Sigma_R-\Sigma_A)P_- - i\Gamma\,\mathbf{1}\Bigr]G_A,
\end{equation}
this can be rewritten exactly as
\begin{equation}\label{eq:GK-spectral-appendix}
G_K(\omega,q)
=
\sgn(\omega)\,G_R(\omega,q)\,(\Sigma_R-\Sigma_A)(\omega,q)\,P_-\,G_A(\omega,q).
\end{equation}
This form makes the large-\(\Gamma\) behavior transparent.  Since
\begin{equation}
G_R(\omega,q)= -\frac{2i}{\Gamma}\,\mathbf{1}+O(\Gamma^{-2}),
\qquad
G_A(\omega,q)= \frac{2i}{\Gamma}\,\mathbf{1}+O(\Gamma^{-2}),
\end{equation}
we obtain
\begin{equation}\label{eq:GK-largeGamma-appendix}
G_K(\omega,q)
=
\frac{4}{\Gamma^2}\,\sgn(\omega)\,(\Sigma_R-\Sigma_A)(\omega,q)\,P_-
+O(\Gamma^{-3}).
\end{equation}
Therefore
\begin{equation}\label{eq:PiKK-largeGamma-appendix}
\Pi_{KK}(\Omega,k)
=
\frac{16}{\Gamma^4}
\int\frac{d\omega\,dq}{(2\pi)^2}\,
\sgn(\omega+\Omega)\sgn(\omega)\,
\bigl[\Sigma_R-\Sigma_A\bigr](\omega+\Omega,q+k)\,
\bigl[\Sigma_R-\Sigma_A\bigr](\omega,q)
+O(\Gamma^{-5}).
\end{equation}

At \(r=1\), the bulk continuum starts at \(|\omega|=1\), so
\(\Sigma_R-\Sigma_A\) has no support near \(\omega=0\).  Hence the full KK
bubble is smooth near \((\Omega,k)=(0,0)\).  The singular shell found in the
strict pole-only, zero-temperature approximation is therefore absent once the
continuum part of the self-energy is kept.

This conclusion is confirmed numerically.  Using direct midpoint quadrature with
\begin{equation}
N_\omega=8192,\qquad N_q=512,\qquad \omega_{\max}=640,
\end{equation}
we find
\begin{equation}
\Re \Pi_{KK}(0,0)
\approx
-8.38\times 10^{-5},\quad
-8.08\times 10^{-6},\quad
-6.70\times 10^{-7}
\qquad
\text{for}\quad
\Gamma=20,40,80.
\end{equation}
On the window \(0\le \Omega,k\le 0.04\), the largest resolved line-cut
deviations are
\begin{equation}
\max |\Re[\Pi_{KK}(\Omega,0)-\Pi_{KK}(0,0)]|
\approx
5.23\times 10^{-6},\quad
9.62\times 10^{-7},\quad
1.21\times 10^{-7},
\end{equation}
\begin{equation}
\max |\Re[\Pi_{KK}(0,k)-\Pi_{KK}(0,0)]|
\approx
1.24\times 10^{-7},\quad
1.71\times 10^{-8},\quad
2.21\times 10^{-9}.
\end{equation}
The line cuts and two-dimensional low-energy maps are shown in
Figs.~\ref{fig:KK-full-linecuts} and \ref{fig:KK-full-heatmaps}.  They are
smooth at the origin and decay rapidly with \(\Gamma\), consistent with
\eqref{eq:PiKK-largeGamma-appendix}.

\begin{figure}[t]
    \centering
    \includegraphics[width=\linewidth]{output/full_KK_bubble_dyson_sign_linecuts.pdf}
    \caption{Direct numerical line cuts of the full KK bubble for
    \(\Gamma=20,40,80\).  Left:
    \(\Pi_{KK}(\Omega,0)-\Pi_{KK}(0,0)\).  Right:
    \(\Pi_{KK}(0,k)-\Pi_{KK}(0,0)\).}
    \label{fig:KK-full-linecuts}
\end{figure}

\begin{figure}[t]
    \centering
    \includegraphics[width=\linewidth]{output/full_KK_bubble_dyson_sign_heatmaps.pdf}
    \caption{Direct numerical low-energy maps of
    \(\Re[\Pi_{KK}(\Omega,k)-\Pi_{KK}(0,0)]\) for the full Green's function.
    The origin is smooth, and no wedge-like singular support remains.}
    \label{fig:KK-full-heatmaps}
\end{figure}
 
\begin{comment}
\section{Edge monitoring}

\subsection{Quadratic order}

Now we specify to edge monitoring, where we consider only $\alpha=\ket{+}$, and $\{\mathbb{P}_\alpha\}=\{\mathbb{P}_+\}$. With that, the fluctuation of interest becomes a single scalar field $\delta \Sigma=\mathbb{P}_+ \tau^y F$. The fluctuation action takes the form 
  \begin{equation}\label{}
  iS = -\frac{1}{2} \Tr\left(\delta \Sigma\left(W_\Sigma-W_G^{-1}\right)[\delta \Sigma]\right)-\frac{1}{3} \Tr\left(\delta \Sigma G \delta \Sigma G \delta \Sigma G\right)+\dots
\end{equation} For the ansatz of $\delta \Sigma$ specified above, the kernel $W_G$ has eigenvalue $\gamma^2=\Gamma$. The kernel $W_\Sigma$ is given by 
\begin{equation}\label{}
  \Tr(\delta \Sigma W_\Sigma[\delta \Sigma])=\int_{q,\Omega} F(-\Omega,-q) W_\Sigma(\Omega,q) F(\Omega,q)\,,
\end{equation} where 
\begin{equation}\label{}
  W_\Sigma(\Omega,q)=\int\frac{\rd \omega}{2\pi}\int\frac{\rd k}{2\pi} G^+_R(\omega+\Omega,k+q)G^+_A(\omega,k)+G^+_A(\omega+\Omega,k+q)G^+_R(\omega,k)-G^+_K(\omega+\Omega,k+q)G^+_K(\omega,k)
\end{equation} Here $G^+$ means projecting to the $\mathbb{P}_+$ basis, i.e. $G^+=\frac{1}{2}\Tr_\text{orbital}(\mathbb{P}_+ G \mathbb{P}_+)$.

We numerically compute $W_\Sigma(\Omega,q)$ using the relation 
\begin{equation}\label{}
  G^{-1}=G_0^{-1}-\Sigma_M\,.
\end{equation}
\end{comment}




\end{document} 
